% =====================================================================
%  Harald Høffding, HENRI BERGSON'S PHILOSOPHY: CHARACTERIZATION AND
%  CRITIQUE.  English translation of Henri Bergson's Filosofi. Karakteristik
%  og Kritik (Kjøbenhavn og Kristiania: Gyldendalske Boghandel, Nordisk
%  Forlag, 1914).  Translated from the Danish.
%
%  GENERATED by ebook-method/bergsons-filosofi/tr_build.py from the parts in
%  ebook-method/bergsons-filosofi/tr/ (edit the parts, then rebuild).
%  Made from texts/hoeffding/bergsons-filosofi/transcription.tex (the text
%  of record) on 2026-10-27, in the main session, chapter by chapter; the
%  translator's odd-spot list (tr/oddspots.tsv) was read at the image and
%  led to four new PRINTED AS IS marks in the Danish.  [N] in the margin is
%  the page of the 1914 print, as in the transcription; the notes are
%  endnotes, as printed, and keep their labels "N. (p. P.)".
%
%  TERMS (Hans's choices, 2026-10-27): Anskuelse = perception (Anskuen =
%  perceiving), kept apart from Bergson's Intuition = intuition; Forstand =
%  intellect (Bergson's intelligence); Skuen = vision; Eftertanke =
%  reflection; Oplevelse = experience; Sjæleliv = mental life; Fornemmelse =
%  sensation, Sansning = sensing; Emne = material (Høffding's term for the
%  given); Udvikling = evolution where biological or Bergsonian, else
%  development.  Quotations from French, German and Latin stay as Høffding
%  prints them; what he quotes in Danish (Bergson, Spinoza, Goethe's Faust) is
%  translated from his Danish.  Titles of Danish books stay in Danish.
% =====================================================================
\documentclass[12pt,a4paper,openany]{book}
\usepackage[utf8]{inputenc}
\usepackage[T1]{fontenc}
\usepackage{libertinus}
\usepackage{libertinust1math}
\usepackage{textalpha}
\usepackage[english]{babel}
\usepackage[protrusion=true,expansion=true]{microtype}
\usepackage{setspace}
\usepackage[margin=1.3in]{geometry}
\usepackage{fancyhdr}
\usepackage{hyperref}
\hypersetup{pdftitle={Henri Bergson's Philosophy: Characterization and Critique}, pdfauthor={Harald Høffding}, hidelinks}
\setlength{\headheight}{15pt}
\pagestyle{fancy}
\fancyhf{}
\fancyhead[LE,RO]{\thepage}
\fancyhead[C]{\textit{Høffding: Henri Bergson's Philosophy (1914)}}
\setstretch{1.3}
\usepackage{marginnote}
\newcommand{\opage}[1]{\marginnote{\footnotesize\textit{[#1]}}}
\newcommand{\chapterhead}[1]{\clearpage\phantomsection\addcontentsline{toc}{chapter}{#1}\vspace*{3cm}\begin{center}{\large #1}\\[6pt]\rule{1.2cm}{0.4pt}\end{center}\medskip}
\newcommand{\initial}[2]{{\Large #1}\textsc{#2}}
\newcommand{\chapterend}{\par\begin{center}\rule{2.5cm}{0.4pt}\end{center}}
\newcommand{\nm}[1]{\unskip\textsuperscript{#1}}
\newcommand{\noterhead}[1]{\phantomsection\addcontentsline{toc}{chapter}{#1}\vspace*{6cm}\begin{center}{\Large #1}\end{center}}
\newcommand{\notelabel}[1]{\noindent #1}
\newcommand{\ix}[3]{\par\noindent\makebox[2em][l]{#1}#2\dotfill #3}

\begin{document}

\frontmatter
\pagestyle{empty}
% --- title page ---
\begin{titlepage}
\centering
{\large HARALD HØFFDING}\par
\vspace{2cm}
{\LARGE HENRI BERGSON'S\\[4pt] PHILOSOPHY}\par
\vspace{1cm}
{\large CHARACTERIZATION\\ AND CRITIQUE}\par
\vfill
{\small COPENHAGEN AND KRISTIANIA}\par
\medskip
{GYLDENDALSKE BOGHANDEL\\ NORDISK FORLAG}\par
\medskip
1914\par
\vspace{2cm}
{\small Translated from the Danish}\par
\end{titlepage}

% --- verso of the half-title in the print: the author's works, titles as printed ---
% read whole at the image, blind (reads/pages-read.txt, PDF 9); set inside a dotted frame in the print, the titles in capitals
\begin{center}
{\large HARALD HØFFDING}
\end{center}
\medskip
\begin{footnotesize}
\setlength{\parindent}{0pt}\setlength{\parskip}{1pt}\setstretch{1.1}
\hangindent=1.5em\hangafter=1 DEN ANTIKE OPFATTELSE AF MENNESKETS VILLIE. 1870.\par
\hangindent=1.5em\hangafter=1 PHILOSOPHIEN I TYSKLAND EFTER HEGEL. 1872.\par
\hangindent=1.5em\hangafter=1 DEN ENGELSKE PHILOSOPHI I VOR TID. 1874.\par
\hangindent=1.5em\hangafter=1 OM GRUNDLAGET FOR DEN HUMANE ETHIK. 1876.\par
\hangindent=1.5em\hangafter=1 PSYKOLOGI i Omrids paa Grundlag af Erfaring (1882). Sjette paa ny gennemsete Udgave. 1911.\par
\hangindent=1.5em\hangafter=1 FORMEL LOGIK. Til Brug ved Forelæsninger. (1884). Sjette Udgave. 1913.\par
\hangindent=1.5em\hangafter=1 ETIK. En Fremstilling af de etiske Principer og deres Anvendelse paa de vigtigste Livsforhold. (1887). Tredie paa ny gennemsete og delvis ændrede Udgave. 1905.\par
\hangindent=1.5em\hangafter=1 PSYKOLOGISKE UNDERSØGELSER. 1889. (Videnskabernes Selskabs Skrifter).\par
\hangindent=1.5em\hangafter=1 ETISKE UNDERSØGELSER. Om Muligheden af en filosofisk Etik. --- Om Velfærdsprincipet. --- Forholdsloven i Etiken. 1891.\par
\hangindent=1.5em\hangafter=1 SØREN KIERKEGAARD SOM FILOSOF. 1892.\par
\hangindent=1.5em\hangafter=1 KONTINUITETEN I KANTS FILOSOFISKE UDVIKLINGSGANG. 1893. (Videnskabernes Selskabs Skrifter).\par
\hangindent=1.5em\hangafter=1 DEN NYERE FILOSOFIS HISTORIE. En Fremstilling af Filosofiens Historie fra Renaissancens Slutning til vore Dage. To Bind. (1894--95). Anden gennemsete Udgave. 1904.\par
% "Slutning" is partly covered by ink strokes in the Minnesota copy (zoom z5); the reading follows the book's known subtitle
\hangindent=1.5em\hangafter=1 JEAN JACQUES ROUSSEAU og hans Filosofi. (1896). Anden Udgave. 1912.\par
\hangindent=1.5em\hangafter=1 KORT OVERSIGT OVER DEN NYERE FILOSOFIS HISTORIE. (1898). Fjerde paany omarbejdede Udgave. 1910.\par
\hangindent=1.5em\hangafter=1 MINDRE ARBEJDER. Første Række. 1899.\par
\hangindent=1.5em\hangafter=1 DET PSYKOLOGISKE GRUNDLAG FOR LOGISKE DOMME. 1899. (Videnskabernes Selskabs Skrifter).\par
\hangindent=1.5em\hangafter=1 VED AARHUNDREDESKIFTET. Tale holdt paa Københavns Universitet. 1901.\par
\hangindent=1.5em\hangafter=1 RELIGIONSFILOSOFI. (1901). Andet Oplag. 1906.\par
\hangindent=1.5em\hangafter=1 FILOSOFISKE PROBLEMER. 1902. (Universitetsprogram).\par
\hangindent=1.5em\hangafter=1 OM NOGLE RELIGIONSFILOSOFISKE ARBEJDER fra den nyeste Tid. 1903. (Universitetsprogram).\par
\hangindent=1.5em\hangafter=1 MODERNE FILOSOFER. 1904.\par
\hangindent=1.5em\hangafter=1 MINDRE ARBEJDER. Anden Række. 1905.\par
\hangindent=1.5em\hangafter=1 DANSKE FILOSOFER. 1909.\par
\hangindent=1.5em\hangafter=1 DEN MENNESKELIGE TANKE, dens Former og dens Opgaver. 1910.\par
\hangindent=1.5em\hangafter=1 RELIGION OG VIDENSKAB. 1910.\par
\hangindent=1.5em\hangafter=1 UDVALGTE STYKKER af dansk filosofisk Literatur. Med Indledninger af \textit{Harald Høffding}. 1910.\par
\hangindent=1.5em\hangafter=1 PERSONLIGHETSPRINCIPEN i FILOSOFIN. Stockholm 1911.\par
\hangindent=1.5em\hangafter=1 MINDRE ARBEJDER. Tredie Række. 1913.\par
\end{footnotesize}
\clearpage

% --- the Indhold (Contents) of the print ---
\begin{center}{\large CONTENTS}\\[4pt]\rule{1.2cm}{0.4pt}\end{center}
\medskip
\noindent\hfill{\small Page}
\ix{1.}{The Task of Philosophy}{1}
\ix{2.}{Intuition}{11}
\ix{3.}{Psychology and Physiology}{27}
\ix{4.}{The Philosophy of Evolution}{35}
\ix{5.}{The Psychology of Will and of Laughter}{43}
\ix{6.}{Attempt at Metaphysics}{52}
\par
\begin{center}\rule{2cm}{0.4pt}\end{center}
\clearpage

\mainmatter
\pagestyle{fancy}
% --- p. 1 ---
\opage{1}\chapterhead{THE TASK OF PHILOSOPHY}

1. \initial{I}{n} the last part of the nineteenth century the idea of evolution made its way in the various fields of research and furnished new leading points of view as regards the origin and validity of human knowledge. When Goethe's Faust is not content with the usual translation of the first verse of the Gospel of John, but puts ``the Thought'' (der Sinn) in place of ``the Word,'' and then again has misgivings about this expression too, because it is not thought or meaning that creates all things, and finally translates: ``In the beginning was the Deed,'' --- this expresses the view of the doctrine of evolution. All thought, all meaning that we can find springs from the discipline of reality, through the very actions it compels us to undertake. The validity of science receives its final justification in this, that all our science is a carrying-through of the presuppositions from which men must start if their actions are to lead to the results they desire. Practical necessity lies at the foundation of theoretical necessity and proves its validity. So-called Pragmatism has later, after the discussion of the hypothesis of evolution had died down, carried this line of thought further on purely psychological and historical grounds.

Henri Bergson associates himself with this view, but disputes the conclusions drawn from it. The practical necessity of arranging one's ideas in a certain way can in his eyes by no means prove the validity of
% --- p. 2 ---
\opage{2}the knowledge so attained. On the contrary. Far from practice leading to the true theory, the demands of life, the struggle for existence in all its forms, the constant labor of compelling material things to serve human needs, foster precisely an external, mechanical way of thinking, a cutting-up and scattering of ideas, which makes it impossible to gain insight into the inner coherence of existence.

In the struggle with the material, thought is infected by its object and comes to be dominated by spatial images; it grows used to conceiving all relations as spatial relations and to assuming as external a relation between ideas as there is between the different places and objects in space. Only a wholly external connection then becomes possible. Language itself, the necessity of expressing thoughts in words, leads in the same direction. Every thing and every property is designated by its own word, and one then believes that they are as different from one another as the different words are.

Not only natural science, the doctrine of the material, has been subject to this influence, but philosophy too. Bergson traces it all the way back in Greek philosophy, and it is his conviction --- which I shall discuss more closely later on --- that modern philosophy has not in this respect freed itself from ancient philosophy. The oldest Greek philosophers kept very close to perception, as it presents things immediately; they assumed that it was one of the elements accessible to sensation that passed over into the other elements by conversion or transformation. But with Parmenides of Elea (c.\ 500 B.C.) philosophy entered upon the path it has not left since. The Eleatic philosophers' criticism of the concept of change, and their assertion of the unchangeable as that which alone can be the object of true knowledge, led to exalting the pure concept at the expense of perception, to moving in the world of abstraction, generalization and inference instead of losing oneself in what sensation and immediate consciousness contain. In Plato's doctrine of ideas the pure concepts were established as the innermost essence of existence, in opposition to the constant shifting within the world of sense-perceptions and
% --- p. 3 ---
\opage{3}experiences \nm{1}. And since then, according to Bergson, philosophy has steadily gone in the same direction.

In this, according to Bergson, there is a great misjudgment of the significance of philosophy, which runs through its whole history. This significance consists, he maintains, precisely in freeing thinking, and mental life as a whole, from the predilection for the unchangeable, for what constantly repeats itself, for the fixed and mechanical, which the practical handling of material things leads us to call forth. Science and philosophy have hitherto stood in the service of practical ends and have thereby been led into constantly vain attempts to understand changes on the basis of the unchangeable, the perceptible on the basis of pure concepts. The proper task of philosophy is to try to regain the immediate unity and the inward coherence of our mental life, from which practical life, in whose service language and natural science stand, again and again leads us away. The life of thought must become entirely disinterested in order to be able to solve this task, --- in order to catch sight of what is really given, --- the foundation of what language seeks to express and what science seeks to formulate.

The title of the work with which Bergson first made his name, and which in my eyes is still his best work --- Les Données immédiates de la Conscience (Paris 1889) --- is characteristic in this respect.

2. Henri Bergson was born in Paris in the year 1859, and his most important work thus appeared when he was 30 years old. No detailed information about the course of his life is available; it seems to be at one with the development of his intellectual life. After passing through the École Normale he worked as a teacher at various secondary schools in the provinces and in Paris, then for a time (1898--1900) at the École Normale, and he now works, from the year 1900, as professor at the Collège de France.

He is a brilliant, often dazzling writer and is said to be an excellent speaker. His lectures are overcrowded and seem to have become a fashion among the intelligent bourgeoisie. That is not, however, something he has himself pursued. Personally he makes a quiet, refined and attractive impression, and
% --- p. 4 ---
\opage{4}I for my part remember with pleasure the interesting philosophical conversation I had with him in Paris some years ago. I had already at that time briefly indicated my difference from him in a historical survey, and later I have discussed various important points in his view in my book ``Den menneskelige Tanke'' \nm{2}.

Bergson represents, as will already be seen from the foregoing, a reaction, both against rationalism and against realism, a reaction which can be traced in the present day generally, in various fields. There may be those for whom judgment on Bergson has already been passed with the word reaction. But a reaction can be fully justified when new points of view are needed, or when old points of view, which were for a time pushed aside because research on its winding ways had for a time no use for them, again prove their necessity. The history of thinking shows that it advances through a constant interplay of actions and reactions, which mutually arouse and indirectly supplement one another. Synthesis and analysis, continuity and discontinuity, unity and multiplicity, perceiving and reflection, involuntary grasping and conscious seeking --- all these different directions or tendencies in the life of thought must constantly interact, but in reality this can only happen in such a way that now the one, now the other gains the upper hand. It is now inner, now outer causes that determine which tendency prevails for the time being. Now there is a spiritual need in one direction or another that demands satisfaction, now it is new observations that demand other forms and set other tasks than those that have prevailed until then. My task in what follows will be to examine how matters stand in this % the Danish prints dennne (PRINTED AS IS)
respect with Bergson's contribution to the philosophical discussion.

He conducts his case with a brilliance and an eloquence that have contributed much to arousing attention and assent. And there is moreover a current in the age that meets him halfway; from sides as different as Catholicism and Syndicalism sympathy has been shown him.
% --- p. 5 ---
\opage{5}It may be said at once that the grounding of the thoughts and their exact working-out are not on a level with the brilliance of the exposition. And a main difficulty for a view such as Bergson's presents itself at once. If both language and scientific thought suffer from a radical imperfection, springing from the fact that they have been formed through and through by their work in the service of practical ends, what expressions and what forms of thought are then available for presenting and developing the higher kind of understanding to which Bergson would lead us?

Everyone must admit that reflection must always build, and actually does build, on the foundation of perception, whether by perception one means sensation, memory or imagination. And one might imagine that our power of perception was so comprehensive and so mobile that it could always work in a satisfactory way, without analyzing and abstracting reflection needing to step in. In my book ``Den menneskelige Tanke'' I have mentioned this possibility, but added at once that this is not, after all, how things stand. Conscious comparison and ordering of the materials and elements that make up the content of perception are required again and again, if this content is to appear in full clarity and with its grounds. The blind, involuntary and unconscious way in which materials and elements are combined in immediate perception proves in experience not to be sufficient. Even if one must grant Goethe that one should dwell in perceiving as long as possible and not let dissolving reflection and judging begin their work too early, yet, when mental development does not stand still, the point will come when this work must be done \nm{3}.

Bergson, like Goethe, lays very great weight on perception and regards reflection as a sign of imperfection. ``If our senses and our consciousness had an unlimited range, --- if our faculty of outer and inner observation were unlimited, --- we should never have recourse to the faculty of forming concepts and drawing inferences. Forming concepts is a makeshift in the cases where one cannot observe, and we are compelled to make do with inferences only
% --- p. 6 ---
\opage{6}where it is a matter of filling the gaps in outer or inner observation and extending its range. I do not deny the usefulness of abstract and general concepts --- any more than I dispute the value of bank-notes. But just as a note is only a promise of gold, so a concept has value only through the possible observations it represents. On this, I maintain, all are agreed. For all admit that the most ingeniously built concepts and the most learnedly constructed chains of inference collapse like houses of cards when a single really observed fact contradicts these concepts and inferences. Moreover there is no metaphysician and no theologian who will not be willing to maintain that a perfect being must know everything through immediate perceiving (intuitivement), without needing to pass through the intermediate links of inference, abstraction and generalization'' \nm{4}.

3. Bergson of course admits that we are not, after all, such perfect beings that we can make do with mere perceiving. And even if we were more perfect in this respect than we are, we should still like to communicate to one another what we behold, and that can only happen through language, which according to Bergson has become useless, as science has. Neither concepts nor images suffice when it is a matter of rendering the content of immediate perceptions; concepts stand in an external relation to one another and each expresses only a single side or property of what is perceived, and no image is fully exhaustive. Yet the image has this advantage over the concept, that it gives something individual and concrete, and Bergson therefore invites us to make an experiment with images. Try, he says, to think of yourself, of your own inner life, as it unfolds in time, of your own changing states! Various images may then come forward. There is probably none of them that suffices, but they may all point in the same direction. You may perhaps find a contrast between the outer crust and the inner layers. But this image does not do justice to the constant mobile coherence (la continuité d'écroulement). Or it seems to you as if something were unwinding within you, as a ball of yarn
% --- p. 7 ---
\opage{7}is unrolled, or the reverse, as if something were rolling itself up in you. Neither does this image suffice, because what is unrolled or rolled up is imagined as homogeneous and as extended in space. Try, then, the image of a spectrum with a thousand shades! But these shades too lie outside one another, and moreover they lie fixed; there is lacking the flowing-on that inner observation directly attests. --- No image, then, suffices. But all the images brought forward nevertheless point in the same direction. And precisely in that each image limps, but all images point in the same direction, an impulse can be imparted, a need awakened for the self-active production of an inner state or attitude (attitude) of the kind into which the speaker wants to bring the one he is speaking to. Only indirectly can a wholly immediate experience be called forth. It is only by self-activity, by an act of will, that we come to stand face to face with what is most immediate and involuntary in our inner life, that to which the philosopher would lead us back.

This exposition, % PRINTED AS IS in the Danish: Fremstiling
which is found in the ``Introduction to Metaphysics'' \nm{5}, is perhaps the most interesting thing Bergson has produced. He has seen that every deeper-going spiritual influence can only take place indirectly, and that it therefore takes art to exert it, just as both receptivity and a capacity for self-activity are required in order to be roused to find oneself, to discover what is innermost in one's inner life. Philosophy has known this ever since the time of Socrates, but it has not always made use of it. Both system-building and the forming of schools have been hindrances here.

One must carefully note the conditions Bergson asserts for attaining the immediate vision which for him is the highest, but which the practice of life so easily hides from us. It is by no means, as has sometimes been said, Bergson's intention to favor intellectual laziness. There are said to be young French students who through Bergson's teaching have acquired contempt for scientific work. This work might, they think, be necessary for practical men, for engineers, physicians and others, but for philosophers, whose domain begins where science ends, the thing is to immerse oneself in immediate experiences, in
% --- p. 8 ---
\opage{8}pure vision! --- Bergson has rightly protested against this interpretation of his philosophy \nm{6}. Even where he is thinking of the ``immediately given,'' which must be accessible to all if they will only work themselves free of the ways of seeing of practical life and the dissecting mechanism of language, even there, as already remarked, both receptivity and effort are required in order to be able to return to immediate experiences, just as there are, of course, peculiar difficulties to overcome when one wants to express what one thus experiences. One cannot, then, dream oneself into ``the immediately given.''

Now Bergson does not, as will appear in what follows, distinguish clearly between ``the immediately given'' in the sense mentioned and the intuition, the immediate vision, which is one with a total view that can stand as the conclusion or rounding-off of scientific thought. Insofar as he speaks of an intuition of this latter kind, it goes without saying that it can only be reached through the purgatory of reflection. But there are of course always those who would try to get into paradise without passing through purgatory.

It is, in my view, as the following exposition will try to show, rather a kind of artistic vision that Bergson prepares the way for than a higher science. Such a science, however high it may be, still constantly needs concepts and inferences, images and comparisons. Who vouches for the validity of the intuition that is supposed to be able to arise after one has passed through the purgatory of research? That by the indirect route Bergson describes so well an inner movement, a deepening and a quickening within the world of the mind, can be aroused, is psychologically possible; but wherein lies the reason that something should thereby be able to appear which, as regards validity, as regards connection with what is innermost in existence, stands above the points of view and results of science? ---

4. By his strong emphasis on the immediately given and on immediate living and beholding, in contrast to reflection, its forms and its results, Henri Bergson recalls a series of earlier French writers, without his
% --- p. 9 ---
\opage{9}being able to be said to have been influenced by them. He belongs to a type that has asserted itself again and again in French intellectual life.

I name Montaigne first. Behind his skepticism stands a devotion to nature, ``our great mother,'' who stirs immediately in all of us, in each according to his peculiar character. And nature in us and above us cannot be expressed by means of sharply delimited concepts. --- Pascal sets the immediate experience of ``the heart'' against abstract demonstration, ``esprit de finesse'' against ``esprit de géométrie'': the former sees the thing in its peculiarity all at once and with a single glance, the latter grasps it through a series of inferences. --- An analogy to this is offered by Rousseau's contrast between nature and culture, feeling and intellect, collectedness and dispersion. --- Maine de Biran emphasizes various stages of immediate living: vie animale, vie humaine and vie de l'esprit, and he dwells especially on ``the inner man,'' who cannot express himself outwardly in such a way that full understanding becomes possible. From Maine de Biran a line of development runs through Félix Ravaisson, Lachelier and Boutroux to Bergson. --- Bergson offers an analogy with Guyau as well. Guyau's main idea is the radiation and unfolding (expansion) of the powers of the mind. Thought plays a subordinate role in relation to this. And in his poem ``Le mal du poète'' he depicts the poet's malady as consisting in this, that no image is able to express exhaustively the fullness of experience.

A corresponding significance was had in the history of German philosophy by Jacobi, Herder and Hamann over against Kant at the turn of the nineteenth century. They too asserted the collected, concentrated activity of the spirit in opposition to analysis and criticism, the living, concrete coherence of things in opposition to the sharp distinctions of the critical philosophy. And in a later time Friedrich Nietzsche pleads, in overflowing passion, the cause of the mighty unfolder of life against rationalism \nm{7}. ---

Such intellectual tendencies have this significance for philosophy, that they draw attention to how great a content will always remain after thought has done its best to analyze, distinguish and combine the materials it is
% --- p. 10 ---
\opage{10}occupied with. Thought always draws from a fullness that it will never be able to exhaust. This will not, however, daunt its courage or lessen its significance; on the contrary, it proves that it will never lack tasks.

The greatest figures in French philosophy, Descartes and Comte, do not belong to this type. The last word of philosophy must be the right and authority of thought, which need not entail a misjudgment of the manifoldness of life. In the present, Comte's positivist tendency is represented with energy by the sociologist Durkheim. As representatives of Descartes' rationalist tendency may be named, besides Émile Boutroux, two recently departed thinkers, Alfred Fouillée and the mathematician and philosopher Henri Poincaré. Henri Bergson stands in definite opposition to both tendencies. We shall now examine more closely in what follows from which points of view he opposes them.

\chapterend

% --- p. 11 ---
\chapterhead{INTUITION}

\opage{11}5. \initial{T}{here} is according to Bergson a great contrast between what appears immediately and involuntarily in and for consciousness, and what reflection distinguishes, separates and then afterwards seeks to connect. Prior to and beyond all separating and connecting there goes a constant stream of experiences, of inner changes in unbroken, living coherence, which forms the necessary foundation on which reflection can unfold, and at the same time stands as the ideal toward which we seek back from the work of reflection when the work of thought is to be concluded.

Intuition, by which Bergson means immersion in this immediate, flowing coherence, is for him both the foundation and the conclusion of the work of thought. And I cannot but see that there is a decisive unclarity in him as regards the relation between intuition as a psychological presupposition and intuition as the conclusion of thought, an unclarity which has partly caused the misinterpretations his philosophy has been exposed to.

From his very first appearance Bergson places himself in justified opposition to a tendency which in the preceding generation exerted a great influence in the philosophical, especially the psychological, field. It is the tendency usually designated as English association psychology, and which is marked by the names of Hume, James Mill, Stuart Mill and Herbert Spencer. This tendency has had the great merit of having prepared the way for a more scientific psychology by conceiving mental life as consisting of a multitude of elements --- sensations,
% --- p. 12 ---
\opage{12}ideas, feelings and impulses --- which were independent of one another and could come into firmer connection with one another only in an external way, namely by often occurring together. Mental life was a chaos which only by the power of habit came to form an ordered whole. This theory has contributed much to leading beyond vague descriptions and mystical references to the unity of mental life. It has made it easier to put definite psychological questions. And it has contributed much to bringing psychology into closer connection with natural science. But it decidedly confused its theoretical and abstract points of view with what observation of the mental life actually given shows, and it came especially into conflict with the fact that no mental element can be understood in its peculiar character except through the coherence and interaction in which it stands with other mental elements. The elements as such are always products of analysis, of reflection, and must not be confused with what immediately given mental life presents.

This view was opposed both in and outside England. Here may be mentioned Wundt's ``Physiologische Psychologie'' (1874), with its strong stress on the significance of attention (``apperception'') for the development of mental life. The same year Thomas Hill Green's ``Introduction to Hume'' appeared, with an acute criticism of the presuppositions from which the most famous representative of ``association psychology'' had started. The same year my book ``Den engelske Filosofi i vor Tid'' also appeared, in which the critical remarks already indicated the view that was carried through in my ``Psykologi paa Grundlag af Erfaring'' (1882), especially with the help of the relation between healthy and sick mental life. In England Francis Herbert Bradley's ``Ethical Studies'' (1876) marked the decisive break with the traditional English school \nm{8}. In the field of descriptive psychology William James' essay on the stream of thought (The Stream of Thought. Mind 1884, later reprinted as ch.\ 9 of his ``Principles of Psychology'') made an epoch through a brilliant depiction of the great involuntary coherence of mental life, which is the foundation of all self-observation and all psychological analysis.
% --- p. 13 ---
\opage{13}Bergson is thus not the first who has criticized the transfer of purely mechanical analogies to the field of mental life and sought to let what is actually given come more into its own. But he has done it with a force and an art of exposition that far surpasses his predecessors', in his chief work ``Les Données immédiates de la Conscience'' (1889), whose line of thought is now to be reproduced.

6. Our mental experiences take place --- as we shall discover when we free ourselves from the spatial and mechanical ideas that have forced their way from natural science into the sciences of the mind --- like gliding, coherent streams. They take place in time, not in space. One state passes immediately over into the other, perhaps just when this is of a wholly opposite kind. It is no homogeneous stream; qualitative changes constantly assert themselves. Sensation of color succeeds sensation of color, sorrow succeeds joy, attention dullness, tension rest, the great the small, and so on. And these experiences are each understood only through their coherence. They are not separated from one another in the same way as the things in the external world, each of which is in its own place in space.

A consequence of this is that mental experiences cannot be measured directly, as what takes place in space can. We cannot determine whether two sensations are absolutely equally strong. We cannot add and subtract sensations. When we say that we measure them, the measuring takes place through their relation to the external stimuli; as when we measure sensations of heat by means of the thermometer. The concept of magnitude does not belong in the sciences of the mind at all, but in the science of material phenomena. In the struggle for the physical maintenance of life it is important to be able to count, measure and weigh, and it is the greatness of natural science that it is able to carry through a quantitative way of seeing for the whole material world. But all the work done here is from the first called forth by the practical needs of life and stands constantly in their service.
% --- p. 14 ---
\opage{14}The attempt Fechner has made to measure mental phenomena accordingly becomes, as a consequence of this way of looking at things, the object of criticism on Bergson's part. He keeps here, however, exclusively to the difficulty of finding psychic units, that is, sensations so simple that they can be substituted for one another with absolute exactness, and less simple sensations be regarded as sums. This criticism is not new; it was --- as Bergson himself mentions --- put forward as early as 1875 by the mathematician Jules Tannery \nm{9}, and its correctness is generally recognized. Bergson strikingly expresses its main idea in the following words: ``The most essential, in my opinion the only question is whether a contrast between A and B can really be equal to a contrast between B and C, when the two contrasts take place between different elements. As soon as one has been able to establish that two sensations can be equal without being identical, psychophysics will be founded.''

There are other sides of Fechner's psychological investigations that Bergson leaves aside, although they would fit well into his theory. It is the way in which Fechner gives an experimental proof that a sensation has no absolute magnitude that is the same under all conditions, but that the magnitude or strength with which it appears depends on its (or rather the corresponding stimulus's) relation to preceding or simultaneous stimuli.

Moreover it might have been pointed out that all measurements in the material field too rest on analogy, on a transfer of numerical relations to the relation between spatial phenomena, and that here too it is not possible to point to units that are absolutely unchangeable under all conditions.

7. A significant peculiarity of all mental phenomena, when they appear in their full immediacy, is that they unfold in time. The concept of time accordingly plays a chief part in Bergson's philosophy. Now this concept is of course also an important concept in natural science, which is more and more becoming a doctrine of motions. But it is precisely here that Bergson finds
% --- p. 15 ---
\opage{15}the most fateful point where unwarranted material analogies from natural science have forced their way into the sciences of the mind. In natural science one dissolves --- he thinks --- time into moments, and one then conceives each moment as a place or a small stretch. Two times are called equal in natural science when like bodies under like conditions have traversed equal stretches. One then puts space in the place of time. Or rather, psychological time is replaced by mechanical or geometrical time. In psychological time, the time in which we actually live and in which we experience what is the actual content of our life, every moment has its peculiar character, is filled in its own particular way, and no moment is therefore quite like another. --- Psychological time Bergson calls la durée, more exactly la durée qualité, as distinct from the temps matérialisé of natural science, which can be symbolized by a stretch in space.

The tendency to use spatial images even when the talk is of our inner experiences stems, according to Bergson, not only from the influence of practical work and of natural science, but also from the need for communication. Space affords us the best signs and symbols when we want to make ourselves understood to one another. ``We have an involuntary tendency to make our experiences fixed and solid in order to be able to express them in language. Hence it comes that we confuse our very feeling, which is in a constant becoming, with its external object, which has a constant being, and quite especially with the word that designates this object.'' In space we have a form that is both common to us all and measurable and clearly determinable. In the perception of a homogeneous space one is already on the way to social life. \nm{10} For then one has a common schema by whose help one can order and clarify one's communications. Everything can find its place there, if one is not too particular about the difference between the immediately given and its symbolic expression. If, on the other hand, we want precisely to try to catch sight of what is really given, then we must strive to disregard all quantitative determinations and all possible symbols. ---
% --- p. 16 ---
\opage{16}8. There is in this whole line of thought, which Bergson has presented in a brilliant way, much that is justified, and Bergson has great merit in having set the difference of mental states from other sides of our experience in a clearer light by his penetrating depiction of their peculiarities. Some critical remarks must nevertheless be made before we go further.

Among the quantitative concepts that must be kept from being applied to mental life Bergson also counts the concept of number. For it rests, in his view, on discontinuity between places in space. He thinks that we can count only what can be referred to different places. But one can perfectly well count qualitative differences. I can count my thoughts and my changing moods within a certain short time. Holberg's ``Vægelsindede'' had fifteen minds in a single hour. The concept of number certainly presupposes discontinuity, but not necessarily discontinuity between places in space. According to Bergson's own view the immediately given is by no means homogeneous and unchangeable; it constantly presents new differences, and here, then, is a sufficient foundation for the concept of number. One could very well imagine beings who knew no spatial relations and who could nevertheless have formed the concept of number. In the end it is basically always our own thoughts or acts of attention that we count. \nm{11}

The point touched on here is connected with the fact that Bergson in general makes too great and too external a difference between the immediately given (what I usually call the materials) and psychic work. Nothing is given to us, no material appears to us, without psychic work, whether or not we notice this work. When one sets an immediately given (a material) over against the activity by which it can be changed through analysis and abstraction, it is not thereby excluded that this given itself may partly be due to a preceding psychic activity. It presupposes attention in order to come into being at all. It has gathered itself into a certain whole for us by virtue of the same process that goes on with a greater degree of consciousness when we work scattered observations together into a whole. The different strength with which particular elements in the given
% --- p. 17 ---
\opage{17}appear is due partly to relations of contrast, partly to different feeling-tone. Under the influence of such forces it happens that our sense-images, our memories and our fancies involuntarily change character, so that we perhaps only long afterwards become conscious of these changes. That is why I have said in my book on ``the human thought'' that what in each particular case is called immediately given deserves this name only in relation to a certain definite work, whereas from other points of view and in comparison with other materials it might perhaps have to be designated as a work.

If there were the opposition between material and reflection that Bergson asserts, it would be impossible for reflection to arise at all, and it would moreover become impossible to give any description of the immediately given. It would be as if two persons dwelt in us, one of whom had the materials, the other reflection; they would not understand, or even so much as discover, one another. ---

9. It is for Bergson as if a fall had taken place. Seduced by the serpents that the practice of life hatches, we have turned our backs on the immediately given and given ourselves over to dissecting and abstracting reflection.

If we now ask when this fall took place, Bergson refers us to the book he himself regards as his chief work: \textit{Évolution Créatrice} (1907), and we are thereby led from psychology over to biology.

The fall happened when the intellect (l'intelligence) took over from instinct. Instinct stands nearer to life than the intellect does. It is an immediate continuation of growth, of the process of organization. When the little chick breaks its eggshell with a blow of its beak, it acts from instinct, and yet it does nothing but continue the movement that has carried it through the embryonic state. Instinct springs from the activity of life itself. The intellect arises only where instinct proves insufficient, --- as a means of making tools. It is as external and mechanical as instinct is organic. The intellect moves around its object, seeks to find
% --- p. 18 ---
\opage{18}as many different external points of view on it as possible, but is unable to make itself one with it. It draws the object toward itself instead of penetrating into it. Instinct and intellect stand in the sharpest opposition to each other. Of a multitude of derived instincts the intellect may perhaps be able to give a scientific explanation, but a complete analysis of the nature of instinct it will never be able to give; for that it is too formal. Instinct cannot be expressed in concepts of the intellect. It is continuity and discontinuity that stand over against each other here. There are in the world many things that the intellect alone is able to seek, but which it would not be able to find. Instinct would be able to find them, but would never seek them!

If the energy that leads to the instinctive act could awaken to consciousness, be made inward as knowledge instead of being converted into practice, --- if the potential or slumbering consciousness in instinct could awaken to full clarity, --- then the innermost secrets of life would be revealed. But the demands of life brought with them the development of the intellect, so that means and tools could be procured, and with that the fateful step was taken. \nm{12} ---

This doctrine of Bergson's about the relation between instinct and intellect involuntarily calls forth the remark that there is, after all, no reason to expect greater light on the riddle of life from instinct than from the intellect. If Bergson is right that the intellect is practical, this nevertheless holds in at least as high a degree of instinct. For it is surely practical to go straight to the goal where that can be done, instead of taking unnecessary detours. Instinct and intellect are two means of reaching a goal, and each corresponds to its own definite conditions of life. There is no reason to look for anything especially mystical in instinct. It must surely be assumed that a need stirs in instinct, but precisely because ``consciousness sleeps,'' it is a blind need that can reveal no secrets to us.

Bergson assumes a certain analogy between instinct and intuition. Intuition is ``a disinterested instinct.'' \nm{13} But because intuition can be called an instinct, in that it works involuntarily, instinct cannot conversely be regarded as an intuition, and it is only
% --- p. 19 ---
\opage{19}thereby that it acquires its mystical stamp. In itself it slaves for the practice of life just as much as ``the intellect.'' And if instinct is to be able to become disinterested, to emancipate itself from bondage in the service of practical ends, why should ``the intellect'' not be able to do so too? Is not the work of thought, moreover, also a kind of life which, if it could awaken to consciousness of itself and its laws, could also throw light on the nature and conditions of life? Life stirs at all stages --- in thought as well as in instinct and in the process of organization. And everywhere involuntary working and striving precede deliberation and clear consciousness.

Just as it was not well-founded and justified to set what is ``immediately given'' in the psychological sense in absolute opposition to reflection, just as little, then, is it possible to set instinct in absolute opposition to the intellect and to regard them as standing at very different distances from the innermost nature of life.

10. If it is asked whether there is no redemption after the fall, Bergson answers that the thing is to work one's way out of the constant preoccupation with practical interests that brings with it absorption in the external and material. In his ``Introduction to Metaphysics'' he had already made an appeal in this direction, and in his work of natural philosophy on ``creative evolution'' he repeated this appeal with special regard to the relation between instinct and intellect. It costs an effort of will to follow this appeal. For the practical needs of life, in whose service the intellect stands, have to such a degree determined our way of thinking, the whole natural tendency of our intellect (la pente naturelle de l'intelligence), that we work against the current when we seek to get beyond the external, mechanical way of regarding existence. We must by an act of will break with scientific habits, with the very fundamental demands of thought, if we want to reach that enlargement of our understanding, that vision in which we can come face to face with the innermost nature of existence, the vision to which it is the proper task of philosophy to lead us.

Yet ``the intellect'' and its way of thinking do not stand
% --- p. 20 ---
\opage{20}merely as a hindrance to the ascent to the intuition that is the highest task. If the intellect had not awakened and done its work, evolution would not have got beyond mere instinct, which leads exclusively and immediately to movement in a definite direction. The mobility of the intellect and the orientation in the surrounding world to which it leads develop powers that can now be turned inward and used to reawaken the possibilities of immediate vision which in instinct did not get beyond the slumbering state, because the urge to movement at once made itself felt. Now instinct can be freed from this urge, become wholly disinterested, and it becomes possible with intellectual sympathy to immerse oneself in, and feel oneself one with, the stream of life that moves in us and through the whole of existence, of which we are members. It is a larger world we must be led into, a world from which the intellect, on account of its limited tasks, had to isolate itself. \nm{14}

Of the character and validity of the metaphysics whose possibility Bergson sees here I shall speak later. Here the point was especially to show that access to the intuition in which, according to him, metaphysics properly consists is not easy, and that it requires the purgatory of the work of the intellect as its presupposition.

From an epistemological point of view too he has impressed upon us what here, seen biologically, appeared as an ascent from instinct through intellect to intuition. In a discussion in the French Philosophical Society on the meaning of the word ``intuition'' he put forward the following view:

``Intuition is indeed an original activity of the mind, which cannot be derived from the fragmentary and external knowledge by which our intellect, as it is ordinarily used, applies a series of points of view to things from outside. But one must not fail to recognize that the intuitive way of grasping reality is no longer natural to us in the present state of our thinking. In order to attain it we must for the most part prepare ourselves through a long and careful analysis and make ourselves familiar with everything that concerns the object of our study. This preparation is especially necessary when it is a question
% --- p. 21 ---
\opage{21}of general and complex realities such as life, instinct, evolution: scientific and exact knowledge of the facts is the preliminary condition of the metaphysical intuition that is to penetrate to their principle.'' \nm{15}

11. It is clear that there is a great difference between the intuition one reaches after having gone through the work of experience and analysis, and the ``original or immediate intuition'' which is the immediately given and which presents an undivided continuity that is afterwards cut up through the work of the intellect. \nm{16} It is a point left unclear in Bergson whether it should not be possible to return to this original intuition without going the long way of the intellect, --- whether a human being in the midst of life might not free himself from the interests and habits of practical life and regain the immediate view of himself and of things. The artist can do so according to Bergson, as I shall discuss more closely later; he sees things in their immediate fullness and individuality. But should this not be possible for everyone, since something of what stirs on a grand scale in the artist does, after all, stir in everyone? --- It would have been of great benefit to clarity if Bergson had expressed himself definitely and clearly on the relation between the intuition that forms the necessary presupposition of all psychic work (or rather: is its first result), and the one that is supposed to be able to form the culmination and conclusion of the work of thought as a higher unity of instinct and intellect. Since Bergson uses the word intuition for both kinds, misunderstanding has lain near at hand.

The metaphysical intuition is supposed to be a rebirth of instinct. But it is then a rebirth that wholly changes the character of instinct. Instinct is, as we have seen, every bit as ``practical'' as the intellect, but now it is to be wholly disinterested. And whereas it is characteristic of instinct that it finds without seeking (in contrast to the intellect, which seeks without finding), its finding is now to stand as the fruit of a seeking, although it is only by an act of will, by a violent effort, a leap, that one finally reaches the concluding vision.
% --- p. 22 ---
\opage{22}From the immediate perceiving of motion and change, which we abandon in our analyses and abstractions, this concluding vision differs by a deepening and a widening; with full consciousness the philosopher concentrates himself in the perception of the continuous process, of which he has the best and the nearest example in what he experiences in his own inner self, where past and present immediately join hands.

It is Bergson's endeavor to bring forward the immediate and involuntary states of mental life. Now these may be very different from one another; but through their common character they come for Bergson to stand nearer to one another than they do in reality. \nm{17} Immediate and involuntary perceiving, as it appears in sensing, memory and thinking, --- the involuntary and blind combination of need and capacity that is called instinct, --- and the vision that can form the conclusion of a strict and intricate work of thought, --- these are, in spite of the analogy there is between them, three very different spiritual phenomena, and it leads to questionable consequences when they are not sufficiently distinguished. It is a certain romantic interest, the same that led Rousseau to seek from culture to nature, that leads Bergson to seek out and dwell on those three phenomena, so very different from one another.

12. There appear in the history of thinking different meanings or forms of intuition, and it will be of interest to see whether Bergson's intuition coincides with any of these, or whether it represents, in the opposite case, a new and justified form.

a. In sensing, in memory and in imagination an immediate perceiving prevails. As soon as I open my eyes I get a sense-image that rounds itself off into a certain whole. In memory earlier experiences come up as involuntary images of wholes. And in imagination there takes place a new formation of such images, whose manner of coming into being can be as unconscious as the way in which a sense-image or a memory-image comes into being. There will not from the first be any sharp difference between sensing, memory and imagination, and they will in many
% --- p. 23 ---
\opage{23}ways fuse and combine with one another. It is only critical reflection that leads to a sharp distinction between these three kinds of immediate perceiving, since cases of doubt compel us to seek definite criteria in order to avoid confusing observation, memory and imagination.

Intuition in this sense concerns something individual and special, and stands in clear opposition to the distinguishing and attentive activity we call analysis, as well as to abstraction and generalization and to the drawing of inferences.

b. Intuition can also mean an estimate, an involuntary decision. Here intuition means not immediate perception of an individual whole, as in sensing, memory and thinking, but an assurance or conviction that comes into being through what has been called ``natural suggestion'' or ``natural magic.'' So when certain principles seem to follow ``of themselves,'' or when one is immediately assured of the reality of what one sees before one, or harbors an involuntary trust in an authority. Intuition in this sense of the word, which is very frequent in everyday speech, was set by Thomas Reid against Hume's criticism, and by Jacobi and Hamann against Spinoza's and Kant's philosophy. In the forming of men's view of life and faith this kind of intuition plays an essential role. It is an involuntary synthesis of experiences and observations.

While Bergson lays great weight on the first kind of intuition (a), he does not go especially into this second kind of intuition (b), although it, no less than instinct and intellect, asserts itself in practical life, and although it, no less than the kinds of intuition Bergson describes, serves to lead beyond the division that reflection can cause. And at bottom his own ``intuition'' comes very close to this kind.

c. A third kind of intuition we have in the immediate apprehension of difference or identity between two sense-, memory- or fancy-images appearing simultaneously or successively. It is an act of thought that works in all comparison, and that leads to two sensations or ideas being either kept apart from one another or regarded as one, so that they
% --- p. 24 ---
\opage{24}can henceforth be substituted for one another. Intuition here means immediate knowledge of a relation. It is the single step of reflection, which by being repeated and continued can lead to coherent understanding and survey.

Intuition of this kind evidently does not stand in opposition to intellect, but designates precisely a form of the activity of the intellect. And it stands in opposition to intuition of the first kind (a) only insofar as it concerns the relation between images, not the images themselves.

This kind of intuition is described by Descartes (in his ``Règles pour la direction de l'esprit'') and later by Henri Poincaré (in ``Science et Méthode''). It designates in them the very transition from perceiving to analysis or to proof, a transition which itself cannot, for the time being, be the object of justification. \nm{18}

This kind of intuition (c) does not occur in Bergson, and this partly explains the opposition he asserts between perception and intellect, an opposition that has not without right been called an unfortunate event in the history of thinking.

d. A fourth kind of intuition is the immediate vision of a coherence or a whole, which can be won by attentive and thinking passage through a series or group of members or parts, with separate apprehension of their mutual relations. Thus the course of an intricate demonstration, or the coherence between different points of view on one and the same subject, can finally stand before us as the object of a single glance. Or a personality, whose particular qualities and actions we have observed and thought over one by one, can finally stand before us in its whole individual peculiarity. The more the third kind of intuition (c) has taken part in the discursive passage, the more clearly and distinctly will the whole later be able to stand before us. The work of thought itself (analysis and demonstration) then benefits the perception of the whole, and the two kinds of intuition, the third (c) and the fourth (d), are by no means opposites. The one can prepare the way for the other. It is the lawful coherence found by way of thought that
% --- p. 25 ---
\opage{25}conditions insight into the whole. Both in Descartes and in Poincaré the expression intuition is accordingly used not only of the single transitions of thought, but also, and in close connection with this, of the survey by which a whole of thought can finally stand before consciousness.

Spinoza gave this kind of intuition its place as the culmination of human knowledge. It differs both from mere observation of the particular materials and from abstract knowledge of the laws of their occurrence. It consists in seeing the particular in its individual peculiarity and yet as an interplay, a manifoldness of general laws, through which it stands in inner connection with the rest of existence. It is for Spinoza the perfection of knowledge, but he declares that it is only very little that he has been able to understand in this way. \nm{19} By his description of this kind of intuition he has rather set a great, in the end insoluble, task than himself given directions for its solution. And such a union of the individual and the general as he demands lies in any case on the border between science and art. ---

Bergson's ``metaphysical intuition'' may recall this fourth kind of intuition, but it does not stress as definitely as this the coherence between the activity of thought and the concluding vision. The metaphysical intuition is to come into being by deepening and widening (creuser et élargir) sense-perception (la perception), without any definite method being indicated for this deepening and widening. The particular kinds and transitions of the work of thought become of no importance; the rational element is pushed aside --- and therefore Bergson's metaphysical intuition comes closer to the second (b) than to the fourth (d). In romantic fashion, moreover, intuition is regarded as a regaining of something lost, a regaining that takes place more in spite of the intellect than by its help. In contrast to this it must be maintained that insofar as intuition of the fourth kind (d) has scientific significance, it points forward, not back. It becomes possible through the working-over by thought of intuitions of the first kind (a), through the analyses and deductions for which intuitions of the third kind (c) prepare the way. The richer intuition a is, and the more strictly
% --- p. 26 ---
\opage{26}intuition c is carried through, the more significant intuition d can become. If philosophy could reach such a conclusion as was assumed in the time of the great systems, it would have to be through an intuition won in this way that the goal could be reached.

While Bergson's psychological intuition (of kind a) is of decided importance and has led him to say a word in season, he has not succeeded in giving a clear concept of the metaphysical intuition whose possibility he asserts, or in establishing its philosophical validity and value.

\chapterend

% --- p. 27 ---
\chapterhead{PSYCHOLOGY AND PHYSIOLOGY}

\opage{27}13. \initial{B}{ergson's} philosophy, especially his doctrine of intuition, is to a great extent directed against the influence of natural science on the sciences of the mind. But then a special interest attaches to the way in which the relation between physiology and psychology, or between body and soul, is conceived in this philosophy. For here the relation between the two great divisions of human science comes forward most definitely and most sharply. In his work ``Matière et Memoire'' % PRINTED AS IS in the Danish: Memoire
(1897) he discusses this question.

14. Bergson begins by taking up a strictly scientific standpoint. He will keep to what objective observation shows and what the method of natural science establishes. Everything in nature consists in motion. What happens in the surrounding world and in my own body, brain and nervous system included, is only motion of different kinds and degrees. From the surrounding world motion propagates itself into our body. It reaches through sense organs and nerve fibers to the spinal cord and brain and is continued in outgoing movements. Our body is an instrument for receiving motion from outside and for releasing motion outward. There is in this respect no difference between the brain and the other parts of the body. Only the brain has, more than the other organs, the capacity to preserve received impressions in such a way that they do not at once lead to outgoing movements, but can come to take part in later movements. The body, the brain in%
% --- p. 28 ---
\opage{28}cluded, is an instrument of movement, nothing else. In no degree, in no sense, from no point of view does it serve to prepare, and still less to explain, a representation (représentation). What in our sense-observations (perception) can be explained by the working of the brain is the actions that are begun or prepared or occasioned, but it is not our sense-observations themselves. Still less does the activity of the brain contain any explanation of memory and higher mental activity. And when everything in and outside the brain is motion, it must be added that one must not look for anything in motion but what one observes. The nervous system has only purely physical properties and has no capacity for anything but receiving, inhibiting or continuing motion. We sense only motion --- and motion can produce nothing but motion. The whole effect of the material processes is exhausted in the work of adapting movements (travail d'adaptation motrîce). % PRINTED AS IS in the Danish: motrîce
If the brain is damaged, it is only movements that are inhibited or hindered, nothing else \nm{20}.

Such propositions one might expect to find in a consistent materialist. But Bergson differs from materialism in denying that material motions can bring about mental or conscious life. The question then presents itself with renewed emphasis how, in his view, motion in space, especially in the minute parts of the brain, is related to mental states.

15. To sense and to recognize means for Bergson, in the simplest case, only to notice a difference between what has significance for our need and what has not, and that one notices this difference shows itself in its determining the movements that are released. But this does not take us, he thinks, beyond the material. Sensation is only a part of matter, or there is at any rate only a difference of degree between the sensation of matter and matter itself. It is all the same whether we speak here of sensation (sensation), observation (perception) or recognition (reconnaissance). There is, as a limiting case, an instantaneous recognition of which the body on its own
% --- p. 29 ---
\opage{29}(tout seul) is capable, without any memory proper coming in; it consists in action, not in representation (représentation).

Our body extended in space thus has the capacity both for sensation and for movement. The material side of our existence is a whole of sensation and movement \nm{21}. ---

But --- one involuntarily asks --- if sensation (observation, recognition) is only movement, like everything else in ``the material side of our existence'' (la matérialité de notre existence), why then does Bergson use two terms, why does he distinguish between sensation and movement at all? He seems here to forget his own warning against language, which in this case would have formed one word too many. That besides motion there also exist sensations he cannot know from external observation, which, as he rightly remarks, shows us only motions in space. Here self-observation must step in. It is by virtue of it that we, for example, distinguish between the pain and the writhings of the body in which it expresses itself. Could we observe the brain in all its different states, we should see only motions and rearrangements of atoms, nothing else. That at the same time something is felt or thought we must know by another route.

When Bergson without further ado ascribes to the body the capacity to distinguish (to sense, observe, recognize), he in reality takes back his first proposition that everything in the material world consists in motion and that motion can only produce motion. Or in other words: within what he calls distinguishing or sensation the whole great problem of the relation between soul and body comes up again, in the form of the question: how is sensation in the sense of motion related to sensation as a phenomenon of consciousness (la perception en tant que consciente)? And he must then admit that what the brain can explain, ``as regards our sensations, is only our begun or prepared actions, % the Danish prints Hindringer (hindrances), a misprint for Handlinger: Bergson's ``nos actions commencées ou préparées''
not our sensations themselves \nm{22}.''

So it often goes when one would avoid the great problems. Even if they are reduced to apparently small dimensions, they keep their sting.

16. While according to Bergson there is supposed to be only a difference of
% --- p. 30 ---
\opage{30}degree between sensation and (other) matter, he asserts a difference of kind between sensation and memory.

It is memory that for Bergson establishes continuity, the inner coherence of mental life. We can often observe how memory-images grow forth and increase in clarity and distinctness. Here we stand before the inner stream to which Bergson so often points, and which he is able to describe in so captivating a way. When there often seems to us to be greater coherence between external objects than between our inner states, Bergson maintains that this rests on inexact observation. ``On looking more closely one will discover that our memories form a chain, and that our character, which is always present in our decisions, is the actual synthesis of all our previous states. In this condensed form our past mental life has even a higher degree of existence for us than the external world; of this we always apprehend only a small part, while we make use of the whole sum of our lived experience.'' \nm{23}

Instead of the usual opposition between matter and soul Bergson thus sets the opposition between sensation and memory.

But what does that help, when ``sensation'' according to Bergson is in its simplest form purely material? He does indeed call his book ``Matter and Memory,'' and rightly, if ``sensation'' belongs under the material. The question remains the old one: how can consciousness (memory) arise from material motion or pass over into material motion? ---

And at bottom Bergson himself denies that there are pure, that is, absolute, isolated, independent sensations. Only as a limiting concept, an ideal case, is one supposed to be able to speak of such. They could fill only a single moment; but Bergson denies that time can be dissolved into absolutely single moments. Modern psychology has shown, in Fechner's footsteps, that every sensation is determined not only by the present sense-stimulus but also by those immediately preceding. Already here the past works together in determining
% --- p. 31 ---
\opage{31}the present, so that one may rightly call every sensation an elementary memory. It is not only in our decisions, as Bergson has strikingly described, but in all our experiences, even the most elementary, that the past works along. Bergson's criticism of Fechner has led him too far. Even if Fechner did not succeed in founding a mathematical psychology, the facts he brought forward nevertheless show that outside memory proper too there is a direct joining-together of our experiences. Bergson could, then, have extended his conception of memory as the expression of the coherence of mental life further down, to the most elementary mental functions. And he does in fact admit that in those real, concrete sensations memory takes part: dans la perception concrète la mémoire intervient! --- ``Concrete sense-observation'' is expressly set in opposition to the pure, ideal sensation which would take place in an absolute moment, but which exists ``en droit plutôt qu'en fait!'' \nm{24}

The difference in kind between sensation (sense-observation) and memory must therefore be given up. That a difference can be made at all between sensation (observation) and memory is due to experience, through the disappointments that teach us that one cannot act on all the ideas that come up. Every memory has from the first, like sensation, the tendency to release action. The first form of memory is expectation or impulse. We begin with belief in all our ideas as well as in all our sensations. It is in the school of experience that we learn to refer some ideas to the world of the past once and for all, while other ideas continue to determine our expectations and our ends. \nm{25}

17. When Bergson seeks to hold fast continuity both within mental life and within material nature, in the one case by means of the concept of memory, in the other by means of the concept of motion, it might have lain near at hand for him to join the hypothesis founded by Spinoza, which assumes a relation of proportionality or analogy between consciousness and motion,
% --- p. 32 ---
\opage{32}without, however, consciousness being derivable from motion or motion from consciousness. It is a hypothesis that gives no solution, but which indicates clearly how we must conceive the relation before we enter upon further-reaching speculations. Bergson recognizes it as a working hypothesis, which is in reality a great recognition, since the warrant of a hypothesis must show itself especially in the use one can make of it. And besides, it is in the end an impossibility that a ``working hypothesis'' which proves applicable throughout, in that it sets new tasks and in that the results found agree with it, should have no more significance than a tool one drops when one has no use for it. It is a contribution to the characterization of existence that we can understand it only by quite definite routes and on the basis of certain definite assumptions. An even greater recognition Bergson pays the hypothesis mentioned when he objects against it that it is only ``a disguised statement of the facts.'' Wherein the disguise really consists I do not know; but a presentation of the facts is, after all, not so bad a foundation for further-reaching reflection \nm{26}.

That Bergson does not lay greater weight on such ``statement'' is connected with the fact that natural science is for him a slave in the service of practical life and therefore has no claim to make contributions of principle to our whole world-view. And yet it is clear that the problem of the relation between the spiritual and the material first gets its proper sting from the fact that natural science has developed into a doctrine of material nature as a great, continuous and unsurveyable coherence of processes of motion, --- as Bergson himself, indeed, has clearly seen. He has even contributed much to sharpening the insight that one has no right to assume that what goes on in the brain is anything other than motions, not processes of a more or less mysterious kind. But from this he then draws the unwarranted conclusion that the brain and the organism have no significance whatever for the life of consciousness except precisely in those moments % PRINTED AS IS in the Danish: Øjeblikkke
when it is to intervene in the material world: otherwise consciousness leads its life on its own, and
% --- p. 33 ---
\opage{33}a life that comprises far more possibilities than material nature can realize.

In a recently published essay on soul and body Bergson has expressed himself on this as follows: ``He who could see into the interior of the brain when it is in full activity, follow the coming and going of the atoms and interpret everything they do, would certainly know something of what goes on in the soul; but it would be only very little that he knew. He would know precisely what can be expressed in the gestures, attitudes and movements of the body, --- that in a mental state which consists in an action that is in the course of being accomplished or perhaps merely of being born; everything else would escape him. He would be, in relation to the thoughts and feelings that unfold in the interior of consciousness, in the same situation as a spectator who sees clearly everything the actors do on the stage, but does not understand a word of what they say.''

By this view both psychology and physiology are really violated. --- Psychologically seen, will and action cannot be isolated from the other sides of mental life. An action ``is born'' at every single moment; for there is not a sensation, not a feeling, that does not bring an element of will with it. It is a false abstraction when the life of memory is separated from action in the way Bergson does it. He makes the life of consciousness far too much a dream-life. And even in dreams, even in the ``pious'' wishes into which the original sanguine urge to act can be converted, the constant prompting to action has not wholly vanished. The soul is here, to use an image of Helen Keller's, like a bird that would work its way up against the wind. If the moments of action have physiological parallels, then all other mental states must have them too. And it must then surely be possible for a psychologically expert observer to interpret the brain processes % PRINTED AS IS in the Danish: Hjærneprocesserne
by way of analogy, even if perhaps the moments in which one proceeds to outgoing action would be the most easily understood. --- Physiologically seen, over against the assertion that far more is contained in the soul than the body can express, one must recall a proposition of Spinoza's whose validity still holds good:
% --- p. 34 ---
\opage{34}``No one has yet learned from experience what the body, regarded as mere body, can do according to the laws of nature alone, and what it cannot. For no one yet knows the structure of the body so exactly that he could explain all its functions.'' --- It is precisely the great significance of Spinoza's view on this point that it at once calls on both physiology and psychology to continue working at tracing the inner coherence in their two fields, in order thereby slowly to prepare the way for an understanding of the mutual relation of these two fields. ---

In the proposition that mental life contains more than the body can ever express, the body that is only an instrument for the soul when it is to act, Bergson finds the possibility of the immortality of the soul; but the question whether this immortality holds for an infinite or only for a limited time, a time x, cannot be answered by scientific means. His line of thought here coincides with a hypothesis which Kant (although it was only ``ein bleyernes Waffen'') recommended for polemical use against dogmatic materialists, and which William James has later taken up. The body is then conceived, in Platonic fashion, as a narrow channel into which the stream of mental life is pressed, and which constantly limits it, without being able to stop it and without being able to prevent it from uniting again with the great currents of the sea. \nm{27}

\chapterend

% --- p. 35 ---
\chapterhead{THE PHILOSOPHY OF EVOLUTION}

\opage{35}18. \initial{I}{t} is a sign of the times that every philosopher in our day must take a position on the concept of evolution. This concept has now taken its place among the fundamental forms or categories of thought, in that it contains directions for putting questions and tasks. Of every event and every phenomenon it must be asked through what stages they have arrived at their present character. It is a more special form of the general concept of cause. \nm{28} A philosophy is characterized in an essential respect by the position it gives this concept and by the way in which it applies it.

For Bergson it must have lain especially near at hand to give this concept a prominent place. He maintains, after all, that both the immediately given (that which we face and are ourselves a part of when we work our way out of the dissecting, mechanical way of looking at things that the practice of life forces upon us) and the concluding perception that we reach through the purgatory of the work of thought show us a continuous movement, a progressing stream within which new properties and peculiarities arise and unfold. Everything is for him activity and development; there is no dead being, no fixed things, no absolute rest.

But the scientific conception of the concept of evolution is according to Bergson unsatisfactory. It lays far too great weight on external relations. Evolution comes on this view to stand as a piecing-together of elements that previously existed in isolation, --- as a combination of separately arising beginnings. A real, inner coherence is out of the question here.
% --- p. 36 ---
\opage{36}This external way of looking at things is for Bergson a consequence of the origin and constant task of natural science. It has sprung from the needs of life, and its task is to serve the maintenance and development of life. But its task is not to understand life itself, in whose service it stands. It has been produced by life under certain definite conditions and in order to act on certain definite things; how then should it be able to comprise life itself, of which it is only a single expression and for which it furnishes only a single point of view? It has been deposited on the way in the course of the process of evolution; how then should it be able to find application to this process itself in its entirety?

Spencer used the fact that science arose during evolution through the struggle for existence in countless generations as proof of the validity of the concepts that lie at the foundation of scientific work. But Bergson finds in this very fact an objection to the absolute validity of these concepts. There is, he thinks, no reason to believe that the mode of conception which in the course of evolution has been forced upon us for practical reasons, under the influence of the needs and demands of life, should be the only possible one. Perhaps not all spiritual powers either have been taken up by practical need. To the unused powers one must then turn if one wants to reach a higher, disinterested understanding of life and its laws! Here, according to Bergson, lies the task of philosophy. \nm{29}

But before he attempts to solve this task, he discusses two main forms in which the external and mechanical conception of the evolution of life that he combats has appeared.

19. Life exists in constant struggle with matter. The thing is to handle material things, to master them and transform them. In order to live one needs only the capacity for such a handling of matter; it is not, on the other hand, necessary in order to live that one understand what life itself really is. It is then no wonder that our intellect falls short here. The purely material can be conceived as composed of parts by whose joining-together it comes into being. In a purely material whole we have
% --- p. 37 ---
\opage{37}the parts given first, and the whole is a simple result of their joining-together.

When science now faces a living being, it is natural that it tries to apply to it the same way of looking at things that it is used to applying to everything purely material. Life too is then to be a mere product, arisen by the joining-together of parts. One then seeks to reduce the doctrine of life to physics and chemistry. That is natural enough, when it is merely a matter of exercising such an influence on the living that it can serve as a means for definite ends. One conceives all organization as a kind of fabrication, and every organism as a machine whose parts and whose structure it is a matter of finding. Voilà le point de vue de la science! --- It is in consequence of this point of view that many scientists see it as the ideal to subject the organism to the same mathematical treatment to which the solar system has long since been subjected in astronomy. \nm{30} ---

Thus Bergson depicts the mechanistic conception of life. It follows, in his opinion, straight from the metaphysics from which men naturally start, without being aware that it is metaphysics. Against it he sets a series of considerations that lead in a wholly different direction, in that they are to show that life cannot be reduced to a system of physical and chemical processes.

The mechanistic conception cannot explain the continuity of life, the close coherence between its past and its present. The later stages of the evolution of life are not due to a combination of elements that already existed beforehand and are only now combined, but they can bring something wholly new. New qualities can arise that are conditioned by the past but cannot be derived from it. By the intimate coherence between the old and the new, organic life recalls mental life. In both cases it holds that the different stages of development are incommensurable with one another, and that the series of the stages cannot be reversed. The growth of an organism, and the very evolution of species in nature, presents a definite order of succession which could not equally well be thought reversed. We come back here once again
% --- p. 38 ---
\opage{38}to the difference between the psychological concept of time (durée vraie), according to which all moments are qualitatively different and cannot be exchanged, and the mechanical concept of time, according to which time consists of like moments that can without further ado be substituted for one another, differing only by their number in the series. ---

As untenable as the mechanistic theory, according to Bergson, is the finalist or teleological one, which assumes a preconceived plan that already contains everything that later comes forth in experience in the course of evolution. Here too time in the end becomes useless. Nothing new arises. What experience presents becomes only a repetition of what was already contained in the plan. Finalism is only an inverted mechanism.

This conception too has arisen under the influence of practical ideas. The demands of life make us lay plans which we afterwards seek to put into reality. And we must then keep here to what we can foresee, and that is what can repeat itself when the same circumstances come again. But life is not exhausted by what repeats itself; new properties and forms can arise that cannot be foreseen, although they can perhaps be explained afterwards. Indeed, repetition is possible at all only for a purely abstract consideration. In real time (durée réelle), --- the time whose tooth leaves the mark of its bite on everything, --- no repetition takes place. Everything is changed through and through, and the same concrete reality never repeats itself.

20. In this criticism of the two opposing theories Bergson has a predecessor in Kant's famous line of thought in the ``Kritik der Urtheilskraft,'' which, curiously enough, he does not mention, although it has stated the problem precisely in the way in which it basically presents itself to this day. Kant showed that for human knowledge there are only two ways of explaining a whole, namely either to conceive it as the product of pre-existing parts, or to conceive it as a work produced according to a previously given plan, --- but that neither of these two ways could be applied to the whole that an organism is. While Kant concluded from this only that we must not confuse our way of
% --- p. 39 ---
\opage{39}knowing with nature's way of acting, Bergson thinks he can give a positive solution of the question. It is especially on this point that he attempts to get beyond what the criticism founded by Kant can achieve.

The intellect, knowledge --- so runs his line of thought --- is only a section, a narrowing within a more comprehensive whole. Life goes beyond the reach of the intellect. By its practical way of seeing, the intellect has shut itself out from the great coherence of life. And it is to this great coherence that we must go back. But how can we reach it? Only by living it through! We notice that around the clear kernel the intellect is able to find there is an indefinite fringe, which is fading into darkness and which is of no use in the practice of life, but which points to a greater depth than mechanism and finalism know. \nm{31}

This vague intuition now gains in clarity, according to Bergson, through a closer consideration of the course of evolution of organic life.

In the great struggle with matter life unfolds in different directions and issues in divergent types, which stand as the end points of these directions. Thus the contrast between the plant type and the animal type appears: the former fixed to its place and without consciousness, the latter mobile and with a capacity for consciousness. They have developed from common forms. Within the animal world there are again different lines of evolution, of which the most remarkable are the one that unfolds in the insect type and the one that leads through the series of vertebrates up to the human type. Of these two types the one is characterized by the development of instinct, the other by that of the intellect. Instinct is a simple continuation of growth, a capacity to make use of the organs (the organic tools) in a definite way; the intellect is the capacity to produce mechanical tools. Mechanical invention is the essential function of the intellect. To this day social life is essentially determined by the fabrication and use of artificial tools, and the inventions that as milestones mark the way of progress have also conditioned its direction. Instinct and intellect are two different but equally elegant solutions of one and the same problem.
% --- p. 40 ---
\opage{40}In the vegetating of plant life, in the involuntary working of instinct and in the conscious fabricating of the intellect, the process of life has unfolded along three different paths. Ever since the time of Aristotle the mistake has been made again and again of regarding these three types of life as three successive stages within one and the same line of evolution. But between them there is not a difference of degree but a difference of kind. They are three mutually divergent directions of an activity that had to split in order to be able to grow. \nm{32}

It is a peculiarity of life in general that it is able to reach the same result by very different routes. Thus reproduction is made possible in very different ways within the different types. The power of sight offers an especially clear example of this peculiarity, when one compares the different forms of development of the eye in different animals. \nm{33}

All this bears witness to a striving, a yearning, an original vital impulse or vital upsurge (élan de vie), which makes its way through the world of matter, and which in the manifold types, as well as in the single individuals each by itself, deposits the results it has won in spite of the obstacles against which it constantly breaks. The history of life is the unfolding of an original store of force, an outflowing from a spring whose waters seek to hollow out riverbeds for themselves in the masses of matter. Matter in itself has, as the so-called entropy shows, a constant tendency to fall back into absolute equilibrium. It is against this tendency that life constantly struggles. The evolution of life is not creative in the absolute sense; it is always conditioned by the relation to the material forces and masses. But something expresses itself in it that neither mechanism nor teleology lets come fully into its own, --- a fruitful unity, an infinite fullness. \nm{34} ---

21. As one sees, Bergson starts from an absolute opposition between the organic and the inorganic. However great the problems that constantly assert themselves here, it is nevertheless a great question whether a philosopher does well to take the starting point of his whole view in such an opposition. It is a fact that however much science must constantly
% --- p. 41 ---
\opage{41}hold fast to the difference between life and the inorganic, this nevertheless does not condition any difference in scientific method. One applies the same methods to the living as to the lifeless, and only insofar as one does so does one reach tenable results. And it is a not uncommon view, which is expressed for example by Jacques Loeb \nm{35}, that ``only technical difficulties are to blame for the fact that the artificial production of living matter has not yet succeeded.''

But even if one grants the sharp opposition between organic and inorganic, Bergson's conception is not more strictly grounded than the two conceptions he rejected. The situation is this, that in face of a subject such as organic life one must always operate with analogies. It is an analogy that mechanism uses (so long as ``the technical difficulties'' are not overcome) when it conceives the organism as a machine. It is an analogy that teleology uses when it explains vital phenomena as the fruit of intention and plan. But it is also an analogy that Bergson uses when he conceives life as an ``upsurge,'' a striving, a yearning. He differs from teleology (with which he has, moreover, a certain sympathy) only by taking his analogy from involuntary, half-unconscious mental life instead of from clearly conscious, rationally calculating mental life. Life is for him really of a psychological nature. His doctrine is a kind of vitalism, built on the assumption of a peculiar vital force that contends with the physical forces of nature, an assumption that is more an expression of admiration and wonder at the mighty and rich unfolding of life than a direction for the exact explanation of the special phenomena of life. And when he interprets this vital force as akin to the mental life that we can immerse ourselves in within ourselves \nm{36}, he interprets nature, like the romantic philosophy of nature, from above downward. Life is a struggle with matter, a struggle which, insofar as it succeeds, leads step by step upward toward the world of the spirit. That is poetry, not science, and at bottom, according to Bergson, it is not meant to be the latter either. It is on a bold analogy that he builds here, just as --- according to what will appear later --- his world-view as a whole builds on an analogy that cannot find confirmation in experience.
% --- p. 42 ---
\opage{42}From Schelling's romantic philosophy of nature Bergson's conception of nature differs chiefly on two points. In the first place he leaves inorganic nature wholly to the mechanical conception. It is purely material, and its chief law is Carnot's principle of the constant tendency toward the dissipation of energy and thereby toward the passage into equilibrium. It is with this tendency that the striving of life constantly has to contend. In the second place he assumes a real development through time, whereas the romantic idea of evolution was wholly formal or ideal, really meant only a systematic coherence among the forms of nature, not an arising of the different forms progressing in time under certain definite conditions. \nm{37} In the interval between Schelling and Bergson the concept of evolution has acquired its realistic character, which can no longer be denied. ---

In its most characteristic lines of evolution the process of life has led, according to Bergson, partly to pure instinct, partly to pure intellect. But for him life is not exhausted in these forms. It is possible to live one's way into the deep coherence of which both blind instinct and the analyzing intellect are only sections. When we give up the practical motives that have guided evolution, and cultivate the disinterested sense for the rich reality of life, then that intuition becomes possible which for Bergson is the culmination of mental life, and which I have already discussed provisionally.

Before I undertake a concluding discussion of Bergson's concept of intuition, I insert an examination of two more special psychological questions --- on the will and on laughter --- which will each in its own way throw light on his main ideas.

\chapterend

% --- p. 43 ---
\chapterhead{THE PSYCHOLOGY OF WILL AND OF LAUGHTER.}

\opage{43}22. \initial{B}{ergson's} chief significance lies in the field of descriptive psychology. With the great art of exposition at his command he seeks to lead us beyond the distinctions that practical and theoretical reasons lead us to make, beyond all analysis that tears apart the web of mental life, back to the undivided coherence of our inner life, and thereby to the spring from which we, without always being aware of it, constantly draw when we think about psychological questions. Each single cup that is filled from this spring holds only a small part of what wells forth, and while we draw, the stream constantly goes on. It would be impossible to reproduce the whole stream out of the contents of these cups, and it turns out besides that the contents of two cups are never the same; they have been taken, after all, at different times and at different stages of the stream's course. This image, though it is not taken from Bergson himself, expresses a main point of his philosophy, the assertion of a deep-lying continuity in mental life, a continuity that must not be confused with sameness, since something new constantly comes forth, though in immediate coherence with and in inner relation to what went before.

Bergson's psychological view is, as is seen from this, directed against two sides: partly against the cutting-up of mental life into isolated, mechanically separated elements, partly against the influences of repetition and habit. He does not deny that isolations and dissolutions can take place in mental life, nor that crusts or fixed deposits can form. But he maintains that
% --- p. 44 ---
\opage{44}both are only surface phenomena, and that, as long as life lasts, there is a constant, coherent movement further in.

In two special psychological investigations Bergson's general view comes forward in an especially characteristic way, and it will therefore be of interest to bring them forward.

23. The concept of will is used in modern psychology in somewhat different senses. For some it is a mystical force, like ``vital force'' in physiology, and they consider it most correct to exclude it from psychology. For others it designates certain very complex and derived mental phenomena, will in their usage being used only of resolutions or decisions that presuppose deliberation among different possibilities; on this view too it has no rightful place among the psychic elements. According to a third view, which I for my part hold \nm{38}, the concept of will comprises the whole series of phenomena that begins with the very first expressions of mental life and leads through reflex movements and instinctive actions up to impulses, resolutions and decisions. The direct or indirect relation to action of one kind and degree or another is, on this view, lacking in no mental state, and all sensations, ideas and feelings receive their particular stamp in the individual cases through this relation, and there is therefore an element of will present in every mental state. --- Bergson's view coincides most nearly with this third way of seeing, only that it is somewhat obscured in him by the sharp opposition in which he sets memory, which for him is mental life proper, to action, which brings mental life into contact with matter.

It is the question of the ``freedom'' of the will that leads Bergson into the psychology of the will. He maintains that every decision or resolution that is rightly to be called free, in contrast to what is due to overwhelming and compelling influences and impressions or to purely momentary impulses, must stand in intimate coherence with the whole of mental life. However different the single directions of will may be from one another, they are nevertheless members of one and the same course of
% --- p. 45 ---
\opage{45}development, different stages in one history. Our experiences are, after all, continuous throughout, pass over into one another through many transitions, and it is purely artificial differences we make when for convenience' sake we sharply separate them. The history of a personality is not known by keeping to external habits or detached utterances. The superficial observer often sees only an external, more or less cracked crust. Even where such a crust has formed, it can be broken through in the energetic decisions of will in which the real self (le moi fondamental) expresses itself, this real self that so often remains hidden, but which conditions the coherence of our personality. Only that action should be called free which has this self as its cause. In such actions the whole previous history of mental life takes part; then above all the innermost self (le moi d'en bas) \nm{39} rises to the surface and asserts itself in opposition to habits and to accidents. However sudden and surprising such decisions may be, exact observation will nevertheless show that a process of fermentation has taken place right down in the depths, --- that an increasing tension in the life of feeling and thought has gone before, by which the apparently motiveless actions find their explanation. Our character is ourselves, and our actions are free when they spring from it. Free actions in this sense are rare, and freedom has many degrees. Most of our daily actions are due to momentary promptings or fixed habits; the great expenditure of energy is saved for the decisive situations --- and even where our whole personality ought to vibrate, inertia and ease often come to prevail.

Bergson's conception of the concept of ``freedom,'' applied to the will, coincides entirely with the definition Spinoza has given: ``That thing is called free which exists from the necessity of its own nature alone and is determined to act by itself alone.'' Freedom in this sense does not form an opposite to necessity.

Nevertheless Bergson polemicizes both against ``the error of determinism'' and against ``the illusion of its opponents.'' The error of the determinists is supposed to consist in this, that they make an external,
% --- p. 46 ---
\opage{46}mechanical separation between cause and effect and let the single motives appear over against one another as isolated and mutually independent elements. It is English ``association psychology'' he is thinking of here; but his criticism of it does not strike every view that assumes the validity of the law of causality for the life of the will. This validity Bergson himself assumes, and that in his very use of the concept of freedom as an expression of the inner coherence of character. Indeed, he declares that it is dangerous to define the concept of ``freedom,'' since one then cannot avoid becoming a determinist, and he does not see that he has in fact himself defined it by his use of it. --- The ``illusion'' of the indeterminists Bergson finds in their assuming absolute beginnings of acts of will and thereby tearing these loose from the coherence with the whole of mental life. It is the same enemy he combats here as everywhere in his philosophy.

Bergson is of course right that human actions cannot be predicted with the same certainty as, for example, astronomical phenomena. And he connects this, likewise rightly, with the fact that repetition can be shown in far less degree in the spiritual than in the material field. To this comes that we for the most part know and understand the preceding mental states only through their influence on the following ones; especially when we are to apprehend other people's states of mind we are essentially referred to their results. We are worse placed than the novelist, who can know beforehand to what result the development of his characters leads; we must ourselves live through the others' states of mind in order to understand their significance --- and how is that possible, since we and they each have our own separate personalities? Moreover, even if it could succeed, there would be no time for any foreseeing, since we must not shorten but must live through everything, right up to the moment of decision.

26. With regard to a quite different psychological phenomenon too, understanding hangs together for Bergson with his main philosophical idea of the opposition between real inner life and the external forms and habits that practical and so%
% --- p. 47 ---
\opage{47}cial conditions bring with them. It is laughter and the comic that he has made the subject of a peculiar little work (\textit{Le rire. Essai sur la signification du comique.} 1900).

Between nature and ourselves, indeed between our own mental life and ourselves, there lies, as it were, a veil, because we must work for self-preservation and satisfy the general rules of society. We are compelled to lay weight on particular external properties and relations of things and of ourselves. And language here too increases externality and mechanism by forming words for each particular. We do not come face to face with the real, individual character of things and beings. We do not see things themselves, but for the most part content ourselves with reading the labels that are fixed on them. And this holds not only of external objects; it holds also of our own states of mind, which escape us precisely as regards what is most inward, personal, originally experienced.

But fortunately art is able to lift this veil and show us existence in its whole originality. For art is a work of souls that have been able to free themselves from practical need and its effects. It is especially the great tragic poet who lays bare the innermost of mental life to us. By means of his own experiences he is able to read in the inner selves of other men. He need not have directly experienced the moods and actions he depicts. But he has experienced in himself possibilities which, if they reached reality, would lead to such characters as he depicts. Thereby these characters become parts of his own being, and he beholds as in a vision reality in its full and whole individuality.

What is reached directly and positively in tragic art is reached indirectly and negatively in comic art.

The comic poet builds essentially on observation of other men. In himself he would find nothing ridiculous, even if he looked for it; for we are ridiculous only by that side of our personality which escapes our consciousness. And in others it is only the outward behavior, the surface, that he can use. And that is more or less common to several
% --- p. 48 ---
\opage{48}different individuals. If the comic poet penetrated into the innermost mental life of his models, --- if he reached the inner movement that goes on under the crust, though constantly hampered by it, --- it would be all over with comedy. Comedy has its nature and its significance as a reaction against everything habitual, mechanical, impersonal in human life. It is therefore not pure art, does not have the disinterested character of art proper. There is always an ulterior thought, which aims at correcting and enlightening.

It is a social deed that comedy performs. It always has a social presupposition. Laughter is always the laughter of a group. The ridiculous rests on the contrast with the customs and ideas of a larger or smaller society. Now it is certainly society itself, to a great extent, that produces the stiffening, the crust, that becomes the object of laughter. Society, after all, demands that definite external forms and rules be observed by all, and can thereby come to make men impersonal. But on the other hand society cannot afford that men become mere machines and creatures of habit, so that self-activity falls away and new tasks, work in new directions, are hindered. Therefore comedy works in the service of society and is a social instrument. ---

One sees how Bergson's main idea comes back in his theory of laughter as that which lays bare the disproportion between mechanism and personality, between shell and kernel, between the stiff and the undulating.

27. There is yet another factor in the ridiculous that Bergson mentions, but on which, in my view, he does not lay sufficient weight, although it is a main point. It is that the ridiculous presupposes a sudden interruption of a coherence. Laughter as an effect of contrast does not come into its own. Contrast is according to Bergson only a condition that serves to awaken attention. Experience shows, however, that a more or less conscious expectation must go before, which suddenly dissolves into nothing. So the contrast works in a
% --- p. 49 ---
\opage{49}downward direction --- in contrast to the sublime, where the sudden contrast also works, as when the transition takes place in an upward direction, from the small to the great.

Bergson's own theory deals with only a single example of such downward contrast. We expect, when a human being comes toward us, to meet something living, something individual, something new; and when we then meet instead a machine, a dead repetition of the past, his own or others', this disappointment can release laughter.

Some of Bergson's more special examples too point in the same direction. When ceremonies without content or when repetitions have a ridiculous effect, it is because an expectation dissolves into nothing. The best example here is perhaps the way in which likeness between persons can have a ridiculous effect, something Pascal had already drawn attention to, and which Bergson rightly claims for his theory. ``Really living life,'' he says, ``would never repeat itself. Where there is repetition, complete likeness, we always suspect a mechanism working behind the living. Analyze the impression one gets in face of two faces that resemble each other too much; one will find that one comes to think of two copies cast in the same mold, or two impressions of the same seal, or two reproductions of the same cliché, in short of an industrial fabrication. This bending of life in the direction of mechanics is here the true cause of laughter.'' But this example points beyond Bergson's theory. For in such cases it is not this association of ideas that is decisive. What is decisive is the very passage from the one face to the other; it is the expectation of a new human being that is disappointed. Pascal rightly lays weight on the fact that each face by itself need not arouse laughter; it is the very juxtaposition that becomes the cause.

28. In other respects too it turns out that Bergson's theory fits only a single group of cases.

He knows laughter only as humiliating, as an expression of
% --- p. 50 ---
\opage{50}bitterness and scorn. ``He who laughs,'' he says, ``asserts himself in a more or less proud way and tends to regard the other as a marionette whose strings he himself holds. In this self-absorption we should quickly discover a little egoism, and behind the egoism something that is less involuntary and more bitter.'' --- But we do, after all, laugh too at sudden release and relief, even if it is at no one else's expense, if only an obstacle or a pressure reveals itself in its nothingness. And we can laugh in sympathy with the very one we laugh at, perhaps even in sympathy with the very trait or action that calls forth the laughter. We can also laugh at ourselves, and in such cases it then turns out that Bergson is not quite right in the proposition that we are ridiculous only through that in our being which escapes our own attention. Bergson speaks only of the laughter of scorn, not of the laughter of humor.

And when he maintains that laughter is always the laughter of a group, it must be remarked that it need in any case not be an actually given group that is presupposed. Laughter can presuppose an appeal to truth, honesty and justice, even if these ideas can only be thought realized in an ideal society, and it can then be the expression, now more of a scornful, now more of a humorous mood in one who stands alone in his time and yet does not feel alone \nm{40}.

It is therefore a wrong opposition in which comic art comes to stand to tragic art in Bergson's theory. Does not the former give us, just as well as the latter, the reality of life? Life is, as Plato already saw, to be called both a comedy and a tragedy, both at once. And if art is to unveil reality, which is always individual, then one must take reality both in its limitation and in its greatness, both in its externality and in its inwardness. All individual existence comprises both, and therefore a double art is required to lift the veil that can cover it from us. ---

When, for the rest, Bergson teaches that ``art has no other aim than to remove the practically useful symbols, the conven%
% --- p. 51 ---
\opage{51}tionally and socially accepted generalities, in a word, everything that masks reality from us, in order to bring us face to face with reality itself,'' --- the significance he thereby ascribes to art naturally leads one to recall the significance he ascribes to philosophy in opposition to analyzing and abstracting science, and we come to ask about the relation between art and philosophy. We are thereby led back once more to the concept of intuition.

\chapterend

% --- p. 52 ---
\chapterhead{ATTEMPT AT METAPHYSICS}

\opage{52}29. \initial{T}{he division} of labor forces its way everywhere, in philosophy too. The different philosophical problems move further and further apart from one another, and perhaps in the end only a common type will be demonstrable for them, a type whose origin then poses a new problem. Psychology, epistemology and ethics each take their own peculiar starting points, and it then depends on whether the trains of thought from these starting points can meet in harmony. In Bergson's case it is none of the problems peculiar to the disciplines named that stands as the central one, but rather a problem that has receded into the background for the philosophers of the critical and the positivist tendency, not because it could ever fall away, but because it is bound up with peculiar difficulty, and because its treatment presupposes a treatment of the disciplines named. It is the problem of existence (the cosmological or metaphysical problem). Bergson boldly takes it on. Like Faust he is dissatisfied with whatever else can be known, if he cannot solve the riddle of the world. Like Faust he wishes that he ``may grasp and find the force that works in the world innermost within, the final solution of all riddles.'' To reach this goal Faust gave himself over to magic; Bergson invokes intuition. It is also by its help that he seeks to resolve the opposition that his main idea of the opposition between life, especially personal life, and mechanism poses. In intuition we are beyond this opposition.
% --- p. 53 ---
\opage{53}In a lecture at the Congress of Philosophy in Bologna in 1911 Bergson declared: ``A philosopher who deserves this name has never said more than a single thing, and this one thing he has sought to say rather than really got it said. And he has said only a single thing because he has seen only a single point, and even then it was rather a contact than a sight!'' \nm{41} It is surely Bergson's own experience as a thinker that has put these words in his mouth. Everything he has written turns, as we have seen, on one main idea and on the solution of the problem it contains. And one notices through his many turns of phrase and images, in his constant circling around the one point, and in the struggle with language and with concepts, how it costs him effort to hold fast and express what for him is central in the world of thought. The most valuable side of the great influence he exerts on wide circles is connected with this peculiarity of his. It is to be hoped that this side of his philosophy will work more strongly than the consequence, sprung from misunderstanding, that has been drawn from it, that one could now give oneself over to dreaming vision, because scientific work of thought does not lead up to the culmination. Bergson himself, in order to be able to present his main idea and his solution from all sides, has had to work assiduously and strictly; in order to reach his goal he has treated a series of more special questions and thrown light on them. Especially as a descriptive psychologist he has done a good work. But psychology is for him only the foundation from which he rises to the metaphysical intuition.

Now intuition is a vision, or, as Bergson says, a sight, a contact. And what it brings us is, to begin with, only an image of our own innermost life. How then can we derive from this a metaphysics, which is, after all, to be a doctrine, a theory? It is, after all, only a side or a part of existence that intuition can show us, and how do we get from there to a conception that could hold for existence as a whole?

This question Bergson does not put to himself clearly and definitely, and yet it is the very possibility of metaphysics that is at stake. This much, however, is clear, that he does not, like Hegel, go
% --- p. 54 ---
\opage{54}the way of construction, of dialectical derivation. But what other way is there, then, that can lead from the part to the whole? There is hardly any other way than the one to which Leibniz already pointed when he maintained that all metaphysics was founded on analogy. Kant too came to the same result by a more critical procedure. And it is in reality, as I already showed in my book ``Den menneskelige Tanke,'' the way that Bergson too goes.

Wrongly does Bergson reproach the philosophical systems with declaring without further ado that a part of existence represents the whole of existence. That must happen in every system, every metaphysics, and it happens in Bergson's too. His intuition shows us our inner life as a constant unfolding, a striving, a new-creating. What we thus find as the innermost in the reality that we ourselves are becomes a type, a pattern (modèle), after which we imagine all other reality, when we succeed in raising ourselves above the practically necessary but external way of seeing things. All other true reality can then differ only in degree from the one intuition shows us in ourselves. It is everywhere a striving, a yearning, an upsurge (élan) that lies at the foundation. The kernel of existence is not dead being, but becoming, unfolding, willing (tendance). In ``Évolution créatrice'' Bergson has, by his own statement, wanted to show ``that the whole is of the same nature as the self, and that one grasps it by an ever more complete immersion in oneself.'' \nm{42}

Bergson's method is then not merely intuitive; intuition is only the first step, inference by analogy must take the following steps. Analogies can be highly fruitful for research; they can lead to discoveries and to new ways of seeing; but they do not prove \nm{43}, and especially not when they are not to serve to illuminate one field of observation by means of another, but are to explain existence in general, where no confirmation by experience, no re-testing of the validity of the analogy is possible. All religious and metaphysical ideas build in the end on more or less bold analogies; but such ideas then also overstep the boundary between the world of thought and that of poetry.
% --- p. 55 ---
\opage{55}Let us look a little more closely at how Bergson applies his intuitions and analogies.

30. In his conception of nature he sets, as we have seen, life in a sharp opposition to the lifeless. He stresses the historical character of life. The new hangs together continuously with the old; but it can for that very reason quite well enter into unexpected opposition to it, so that prediction becomes an impossibility. The physical and chemical laws contain only conditions and limits for life, do not state the cause of life itself. And since exact natural science applies a purely mathematical mode of treatment more and more, while on the other hand life, when we really penetrate into it, presents an analogy, which is perhaps more than an analogy, with our mental life, --- it is in the end mathematics and psychology that come to stand over against each other in the contest over the world-view. Mathematical natural science rests on the presupposition that the same elements and the same relations repeat themselves again and again; psychology, on the other hand, shows us a growth, a new formation, that may rightly be called a new-creating. The concepts of growth and creation are for Bergson one: l'idée de création se confond avec celle d'accroissement. We know such growth or new-creating directly from our own mental life, but it is the mark of everything living in its constant struggle against the tendency of matter to end in absolute equilibrium.

``When I,'' says Bergson \nm{44}, ``consider the world in which we live, I find that the automatic and strictly determined evolution of this coherent whole is an activity that is undoing itself (de l'action qui se défait), and that the unforeseen forms that life marks out within it, forms that can themselves prolong their existence through unforeseen movements, signify an activity that is growing (de l'action qui se fait). Now I have every reason to believe that the other worlds are analogous to ours, --- that things go on there in the same way. I know that they did not all come into being at the same time, since observation still today shows me nebulae in the course of concentrating. If it is everywhere the same kind of activity that goes on, whether it is undoing itself or growing, then I can quite
% --- p. 56 ---
\opage{56}simply express this probable likeness by speaking of a center from which the worlds well forth, like rockets in an immense display of fireworks, --- only that I do not conceive this center as a thing, but as a continual welling-forth (une continuité de jaillissement). When God is defined in this way, there is nothing finished once and for all in his being; he is unceasing life, activity, freedom. If creation is conceived in this way, it is no mystery; we experience it in ourselves when we act with freedom.'' ``The thing is only to behold for the sake of beholding, and no longer to behold for the sake of acting. Then the absolute reveals itself very near to us, and in a way within us. It is of a psychological, not of a mathematical or logical nature. It lives with us. It lives in a psychological time (il dure), as we do, only in several respects more concentrated, more gathered into itself.''

From what has been grasped in intuition Bergson, as one sees, descends the scale of analogy to material nature, whose activity is a tendency toward equilibrium, and ascends to a world-center as the expression of the most concentrated life. Physics, biology, psychology and theology gather in this comprehensive intuition into one fundamental view. This is the culmination of Bergson's philosophy.

31. I now go back to examine the foundation from which Bergson has reached this culmination. The essential point here is expressed in the charge he brings against all science and all philosophy, modern as well as ancient, that they believe they can explain motion on the basis of rest, change on the basis of the unchangeable, the living on the basis of the lifeless. There has, in his opinion, developed here an intellectual habit from which we can free ourselves only by an act of will, by a leap over into an immediate vision that lets us feel the pulse of life as it beats not only in us but in the whole of existence. Bergson thinks that the Greek Zeno's conception of motion as resulting from unmoved states is still classical and typical for science, in that science must always conceive time as a sum of single indivisible moments, just as it conceives nature as a sum of indivisible atoms.
% --- p. 57 ---
\opage{57}In this, however, he is decidedly wrong. The ancient conception was certainly inclined to set the unchangeable above the changeable and to hold that there could be no true science of the changeable. But modern thinking and science, with Bruno and Galileo, decidedly took their stand in the world of change and set themselves the task of finding laws for motion and development. Even rest came to be conceived as potential or minimal motion. All the laws of nature that modern science has demonstrated are laws of motion, and the great role played by the concept of evolution also points in a wholly different direction from the ancient doctrine of the unchangeable. The elements or fixed forms at which one provisionally stops designate only stations, resting places, not unchangeable abodes. Neither multiplicity nor unity, neither discontinuity nor continuity can express the working of thought and the nature of existence; each of them is only an abstraction. Real work of thought, as well as everything we can surmise about the nature of existence, presents a constant interplay between continuity and discontinuity and constant oscillations between them, without the one of these points of view being derivable from the other.

32. The reason for the weight the intellect lays on unchangeable elements Bergson finds in its standing in the service of practical life. In his conception of human knowledge he is decidedly a pragmatist. \nm{45} But he breaks with pragmatism, and with the scientific way of looking at things altogether, when he passes over to his philosophy proper, which demands intuition and sets immediate perceiving as the highest. It is for him the task of philosophy to lead beyond the discontinuous and divided to a vision.

But how does this become possible when, as Bergson repeatedly declares \nm{46}, one can indeed go from intuition to intellect and analysis, but not the reverse, just as one is supposed to be able to go only from motion to rest, but not the reverse? --- One must be able to go both ways, to make intuition more distinct through analysis and to gather what analysis separates through new intuition. As we have seen above, intuitions of the first, third and fourth
% --- p. 58 ---
\opage{58}kinds (a, c and d) stand in constant interaction with one another. Bergson ends in a dualism between intuition and analysis, which is supposed to be overcome only by the concentrating act of will to which he appeals, and which he seeks to draw forth by his art of exposition.

When we succeed in reaching such an intuition, everything, according to Bergson, acquires life and movement for us; everything dead and unchangeable vanishes; we feel ourselves drawn, lifted, borne. ``We \textit{live} more than before, and this surplus of life brings with it the conviction that the gravest philosophical riddles can be solved, or perhaps even that they ought not to be posed, that they spring from a congealed view of the universe.'' \nm{47}

But a new problem then arises as soon as we wake up from intuition: how is the problem-free state related to the state that is held in the grip of problems? Why is there a difference between perceiving and intellect, between conviction and justification? In the very states of intuition one forgets that spiritual states of a wholly different kind exist; but these will assert themselves again and again, and they do not find their explanation through intuition itself, which, after all, sees its enemy in them!

The difference between the two states, the one in which problems arise and the one in which they give way to immediate or regained perceiving, points back to a unity when one attends to the fact that the nature of the human spirit shows itself in both states. Bergson and pragmatism are wrong when they cannot find the essence of personality in scientific research as well as in immediate vision. The very methods science works out, and the practical ways life compels us to go, are conditioned by the nature of the human spirit and would be incomprehensible without it. Already in purely biological terms this shows itself in the close relation there is between organ and function. Epistemologically and psychologically we are now so placed that we know the organ only through the function. It is through the study of the way in which human research and action appear that we can come to know the nature of the human spirit. We cut ourselves off from such
% --- p. 59 ---
\opage{59}knowledge when we place ourselves beforehand outside the life of action and research, which, after all, also belongs to life \nm{48}.

And in what language should the content of intuition be expressible, when intuition stands in absolute opposition to the intellect, and ``only the intellect has a language''? \nm{49} The intellect, after all, cannot help us, since it has fallen into bondage in the service of practical life. As we have seen (3), Bergson is clear about this point and seeks, through comparisons that each limp but all point in the same direction, to call forth intuition indirectly, to draw forth the effort of will, the concentration, that absorption in vision presupposes. And here a brilliant art of exposition is at his disposal. But not only that. He is right that every deeper-going spiritual influence, in the intellectual field as in all others, properly speaking always takes place indirectly, by rousing to self-activity. And in science too it can be of importance that a total picture, a collected vision of a coherence, comes into being. That is the work of imagination in the service of science. But scientific imagination differs from artistic imagination precisely in this, that it is not an absolute conclusion, but always serves as the foundation and introduction to new reflection, whereas the artistic image has its value, its end, in itself. Therefore Bergson's intuition becomes art, not science or philosophy. He repeatedly cites \nm{50} art as an analogy to philosophical intuition. But it becomes more than an analogy when intuition enters into the sharp opposition to reflection that Bergson assumes. And of a solution of problems there is in any case no question; rather, as already remarked, in the passage to intuition there takes place a passage to a problem-free state.

It is in the end the opposition between art and science that marks the character of the Bergsonian philosophy. It is like an opposition between holiday and weekday, between poetry and prose. But such, after all, is life: it moves through oppositions, and a brave thought looks these oppositions in the eye, precisely because both belong to life. The thing is to find in prose the hidden vein of poetry, and in poetry not
% --- p. 60 ---
\opage{60}to forget that the way will lead back again to the prose of life. There is a poetry of the work of thought, which pragmatism is inclined to deny, and which most often works with an inner animating power that is set free only at the peaks and seeks to find expression. But it opens itself only to the faithful worker.

33. At the points where intuition really unfolds within Bergson's philosophy, not only the boundary between thinking and poetry is overstepped, but also the boundary between philosophy and theology. As we have seen (30), the concept of God is for Bergson one with the spring or the center from which the evolution of the world unceasingly breaks forth. In all fields, in all world-systems, though most directly in the evolution of life, he finds such a source of activity and living. The opposition that otherwise plays so great a role in his philosophy, the opposition between life and matter, recedes for him in such moments of intuition. It then belongs to the problems that prove unwarranted and groundless for the highest intuition (32).

A philosopher should always be careful about using theological expressions. To theological ideas there correspond within philosophy at most problems, and one need only think a little about the concept of God (which one does not do in the moments of intuition) to see that it contains the greatest problem of all, the problem whether the ultimate ground of value, of the good, can be thought absolutely identical with the ultimate ground of the reality that experience shows us. This great problem one has no right to push aside. And it is sharpened quite especially by the fact that the series of values (the series of means and ends) and the coherence of reality (the series of causes and effects) are each incapable of completion. Higher values, causes lying further back, and effects lying further out in the future can always be thought. In face of the problems that constantly pose themselves anew here, every intuition can only mean a provisional resting place. It can only be a faith, --- a faith whose psychological nature and conditions philosophy has to investigate, --- when a harmony between the series of values and that of causes is assumed \nm{51}.
% --- p. 61 ---
\opage{61}Bergson has believed that at the culmination of intuition he found a unity of physics, biology, psychology and theology (30). He thinks he has a right to use the word creation of the constant production of new phenomena in the different fields. But creation means for him the same as growth, and he imagines the evolution of the world on the analogy of the stream's issuing from the spring. And a spring is, after all, itself a stream, --- the first part of the stream that is said to issue from it!

In his reply to a criticism that appeared in a theological journal edited by Jesuits, Bergson nevertheless maintained that his philosophy establishes ``the idea of a free and creating God'' (l'idée d'un dieu créateur et libre). To this it was answered, with full right, that for Catholic theology God is not merely the spring from which the stream wells out; God does not unfold himself into a world, but he produces it by a creation in a wholly different sense of the word from the one in which Bergson takes it \nm{52}.

Philosophy does best, and does the life of the spirit the best service, by going its own ways and treating the religious problem too by its own methods. Then the account can be clearly settled, and it can be seen how far one can get in the border regions by purely human means. The work in this direction must always be taken up anew. A conclusion of the work of thought cannot be justified as long as experience can always bring new phenomena, situations and problems. A ``rounding-off,'' on the other hand, will naturally take place from time to time, a provisional resting place for orientation before further work is done. But such a rounding-off is not one with conclusion. There will always be something outside the circle within which we round off. Here Goethe's words hold: ``Nie geschlossen, oft geründet.''

Bergson's ``intuition'' differs from earlier attempts in the same direction in being a vision of motion, change, life, evolution. And he himself constantly stresses that no single person can lead us to the intuitions in which thought can reach its conclusion; the work of many thinkers and from many sides is required. And moreover, for him the most valuable intuitions are, after all, those that have passed through the purgatory of the
% --- p. 62 ---
\opage{62}intellect, although a leap must then be made in order to get over into intuition.

Every philosopher of importance is both a symptom of tendencies in the age and a factor that determines such tendencies. On account of the artistic element in Bergson's philosophy it perhaps acquires its greatest interest by being a conscious expression of a current in our time that has sprung from dissatisfaction with what rational thinking and experience have so far brought. By his brilliant art of exposition and his untiring brooding he has brought this tendency to light in a form of principle. Even one who thinks that the power and right of thought have not received the place and the recognition due to them can fully admire the energy with which the eternal border dispute between life and science is here set before the age's eyes. Whatever the outcome of the discussion Bergson has given the impetus to, and of the struggle between admiration and rejection that his writings have called forth, a new interest will surely have been awakened in reflection on the great questions with which philosophy at all times is occupied.

I have of course had to apprehend and judge Bergson's philosophy from my own standpoint, but I have sought to give an objective characterization and to let the criticism develop naturally from it.

\chapterend

\clearpage
% --- p. 63 ---
\opage{63}\noterhead{NOTES}
\clearpage
\begin{small}

% --- p. 65 ---
\opage{65}\notelabel{1.} (p. 3.) On Bergson's conception of Greek philosophy and its constant after-effects, see \textit{La Perception du Changement} (Oxford 1911) p. 6, and \textit{Évolution créatrice} (Paris 1907) p. 333--347. --- On Parmenides see my \textit{Mindre Arbejder.} Second Series. p. 141--149, on Plato \textit{Mindre Arbejder.} Third Series. p. 82--101.

\notelabel{2.} (p. 4.) In the autumn of 1905 I went through ``Les Données immédiates'' in the Laboratory with a small circle of students. In the year 1907 I gave a short characterization and critique of Bergson in ``Lehrbuch der Geschichte der neueren Philosophie'' (later translated into English). In ``Den menneskelige Tanke'' (1910), of which a French edition appeared in 1911, I discussed partly his conception of what is immediately given in the psychological sense (p. 8 ff.; 16 f.), partly his conception of the will (p. 292 f.), partly his position on the question of the possibility of a scientific metaphysics (p. 315--318).

\notelabel{3.} (p. 5.) See \textit{Den menneskelige Tanke} p. 56--58 on the relation between perceiving and judging.

\notelabel{4.} (p. 6.) \textit{La Perception du Changement} p. 5--6.

\notelabel{5.} (p. 7.) \textit{Introduction à la Métaphysique.} p. 4--9. (Revue de Métaphysique et de Morale. 1903).

\notelabel{6.} (p. 8.) See the discussion in the French Philosophical Society on 28 Nov. 1907. (Bulletin de la Société Française de Philosophie. VIII p. 18; 21).

\notelabel{7.} (p. 9.) Most of the writers named in the foregoing are discussed in ``Den nyere Filosofis Historie'' (see the index) or in ``Moderne Filosofer.'' \emph{Félix Ravaisson} is best known for his book \textit{La Philosophie en France au XIX\textsuperscript{e} siècle} (Paris 1868), \emph{Lachelier} for his work \textit{Du fonde%
% --- p. 66 ---
\opage{66}ment de l'induction} (1871) and his essay \textit{Psychologie et Métaphysique} (1885). Both of these last works of Lachelier have been republished in one volume of the ``Bibliothèque de Philosophie Contemporaine'' (1896). --- Of \emph{Alfr. Fouillée}, discussed in what follows, the work of especial interest in this connection is the one he published shortly before his death: \textit{La Pensée et les nouvelles écoles anti-intellectuelles} (1910), which is to a great extent directed against Bergson.

\notelabel{8.} (p. 12.) Cf.\ \textit{Moderne Filosofer}, p. 42--46. --- Cf.\ moreover, with regard to the whole question, \emph{Dwelshauvers}: \textit{La synthèse mentale.} Paris 1908.

\notelabel{9.} (p. 14.) Revue scientifique 1875. Tannery only later named himself as the author of this article. Cf.\ \emph{Bergson}: \textit{Les données immédiates.} p. 50 and Th.\ \emph{Ribot}: \textit{La Psychologie Allemande contemporaine} (Paris 1879). p. 200; 204 ff. --- The following quotation from Bergson in the text is found in the work named, p. 42.

\notelabel{10.} (p. 15.) \textit{Les données immédiates.} p. 98; 103.

\notelabel{11.} (p. 16.) See on this point \textit{Den menneskelige Tanke}, p. 189--193. In his article \textit{Some Antecedents of the Philosophy of Bergson} (Mind. October 1913) \emph{Lovejoy} has shown that Bergson, in his overlooking of the difference between the concept of number and the concept of quantity, has a predecessor in Ravaisson. Since Lovejoy assumes that Bergson has in general been influenced by Ravaisson, he thinks he traces through him an influence of the Schellingian philosophy. Ravaisson had been one of Schelling's hearers in Munich.

\notelabel{12.} (p. 18.) \textit{Évolution créatrice}, p. 150 ff.; 179--182. --- I quote a characteristic sentence: Si la conscience qui sommeille dans l'instinct se réveillait, s'il s'intériorisait en connaissance au lieu de s'extérioriser en action, si nous savions l'interroger et s'il pouvait répondre, il nous livrerait les secrets les plus intimes de la vie. (p. 179).

\notelabel{13.} (p. 18.) \textit{Évolution créatrice,} p. 192: C'est à l'intérieur même de la vie que nous conduirait \textit{l'intuition, je veux dire l'instinct devenu désintéressé, conscient de lui-même,} capable de réfléchir sur son object % PRINTED AS IS in the Danish: object (objet)
et de l'élargir indéfiniment.

\notelabel{14.} (p. 20.) A principal passage is \textit{Évolution créatrice} p. 198: La conscience se déterminant en intelligence, c'est-â-dire se concentrant d'abord sur la matière, semble ainsi s'extérioriser par rapport à elle-même; mais, justement parce qu'elle s'adapte aux objets du dehors, elle arrive à circuler au milieu d'eux, à tourner les barrières qu'ils lui opposent, à élargir
% --- p. 67 ---
\opage{67}indéfiniment son domaine. Une fois libérée, elle peut d'ailleurs se replier à l'intérieur, et réveiller les virtualités d'intuition qui sommeillent encore en elle. % PRINTED AS IS in the Danish: c'est-â-dire

\notelabel{15.} (p. 21.) Bulletin de la Société Française de Philosophie. IX. p. 274.

\notelabel{16.} (p. 21.) See e.g.\ \textit{Matière et Mémoire}\textsuperscript{8}. p. 201 f. Ce qu'on appelle ordinairement un \textit{fait}, ce n'est pas la réalité telle qu'elle apparaîtrait à une \textit{intuition immédiate} . . . \textit{L'intuition pure, extérieure ou intérieure}, est celle d'une continuité indivisée. --- \textit{La Perception du Changement.} p. 8; 23. ---

\notelabel{17.} (p. 22.) Interesting contributions to the elucidation of Bergson's concept of intuition are found in two Swedish treatises, the one by Docent \emph{Malte Jacobsson} \textit{(Henri Bergsons Intuitionsfilosofi)}, the other by Professor \emph{Hans Larsson} \textit{(Intuitionsproblemet.} Särskilt med hänsyn till Henri Bergsons Teori).

\notelabel{18.} (p. 24.) \emph{Descartes}: \textit{Règles pour la direction de l'esprit.} Chap.\ III et VII. (Oeuvres éd.\ Adam et Tannéry. Vol.\ X.). --- \emph{Henri Poincaré}: \textit{Science et Méthode.} Paris 1912. p. 47; 134; 157. --- Cf.\ also \emph{Kroman}'s ``steps of perception'' or ``elementary judgments of value'' \textit{(Vor Naturerkendelse.} p. 97 f.). --- It seems to have been intuitions of this kind that led the oldest Greek mathematicians to set up propositions that only later became the object of proof. See \emph{Zeuthen}: \textit{Den græske Geometri før Platon} (Videnskabernes Selskabs Oversigter 1914). --- See \textit{Den nyere Filosofis Historie}\textsuperscript{2} I p. 211 f. \emph{Léon Brunschwieg}: \textit{Les Étapes de la Philosophie Mathématique.} Paris 1912. p. 105 ff.; 451. Brunschwieg describes ``the opposition that has been asserted in our day between intellect and intuition'' as ``un accident malheureux de l'histoire'' (p. 452). % PRINTED AS IS in the Danish: Tannéry; Brunschwieg (Brunschvicg)

\notelabel{19.} (p. 25.) Cf.\ \textit{Den nyere Filosofis Historie}\textsuperscript{2} I p. 305 (and p. 328). --- \emph{Harold Joachim}: \textit{A Study of the Ethics of Spinoza.} Oxford 1901. p. 180--185. --- The expression ``Samfoldighed'' used in the text is due to H. C. Ørsted. See \textit{Danske Filosofer.} p. 51. --- On the intuition that plays a role in the development of a world-view or a metaphysics cf.\ \textit{Den menneskelige Tanke.} p. 306--315. --- What \emph{Fichte} understands by ``intellectual intuition'' is an immediate consciousness of our own acting, our own self-activity. Everyone who ascribes an activity to himself appeals to this intuition. (Zweite Einleitung in die Wissenschaftslehre p. 5). (Cf.\ \textit{Den nyere Filosofis Historie}\textsuperscript{2} II p. 146). It does not coincide with Bergson's intuition, since it is an abstract thought, --- a presupposition that can be \textit{extracted} from immediate experience, but
% --- p. 68 ---
\opage{68}that does not wholly coincide with any immediate experience. The assumption of activity, of a state's springing from our own being and not from external influence, rests on \textit{inference}, not on perceiving, and Kant had therefore rightly denied that there was any ``intellectual intuition.''

\notelabel{20.} (p. 28.) \textit{Matière et Mémoire.}\textsuperscript{8} p. 251, cf.\ p. 4; 8; 66; 101.

\notelabel{21.} (p. 29.) \textit{Matière et Mémoire.}\textsuperscript{3} p. 65; 93; 149 f.

\notelabel{22.} (p. 29.) \textit{Matière et Mémoire.}\textsuperscript{3} p. 65; 251. --- The usage by which the physiological processes that go on in the sense organs are designated as ``perceiving'' also occurs in physiological writers, even when they (like \emph{Bechterew}: \textit{Objektive Psychologie.} Leipzig 1913 p. 68) expressly declare that they mean to disregard all subjective observation. One then speaks, for example, of ``perceiving organs,'' partly peripheral, partly central.

\notelabel{23.} (p. 30.) \textit{Matière et Mémoire.}\textsuperscript{3} p. 158. --- With regard to the foregoing see p. 60; 159 (difference of kind between sensation and memory).

\notelabel{24.} (p. 31.) \textit{Matière et Mémoire.}\textsuperscript{8} p. 244.

\notelabel{25.} (p. 31.) See on this my \textit{Psykologi} V B, 4; D, 1--2. --- \textit{Den menneskelige Tanke} p. 99--108.

\notelabel{26.} (p. 32.) Spinoza's hypothesis as a working hypothesis: Bulletin de la Société Française de Philosophie. I p. 68, --- as ``constatation déguisée du fait'': Matière et Mémoire p. 247.

\notelabel{27.} (p. 34.) \emph{Bergson}: \textit{L'âme et le corps.} (In the volume Le Matérialisme actuel, which contains lectures by various French scholars) p. 24 f.; 45--48. --- \emph{Kant's} statement is found in \textit{Kritik der reinen Vernunft}\textsuperscript{2} p. 806--808. --- \emph{William James}: \textit{Human Immortality} (1899). --- \emph{Spinoza}'s utterance is found in \textit{Ethica} III, 2. Schol. ---

\notelabel{28.} (p. 35.) Cf.\ \textit{Den menneskelige Tanke.} p. 227--237.

\notelabel{29.} (p. 36.) \textit{Évolution créatrice.} Introduction.

\notelabel{30.} (p. 37.) \textit{Évolution créatrice.} p. 22; 101.

\notelabel{31.} (p. 39.) \textit{Évolution créatrice.} p. 49--50.

\notelabel{32.} (p. 40.) \textit{Évolution créatrice.} p. 137--152.

% --- p. 69 ---
\opage{69}\notelabel{33.} (p. 40.) \textit{Évolution créatrice.} p. 59--94.

\notelabel{34.} (p. 40.) \textit{Évolution créatrice.} p. 105--114; 273.

\notelabel{35.} (p. 41.) \emph{Jacques Loeb}: \textit{Das Leben.} 1911. p. 9.

\notelabel{36.} (p. 41.) \textit{Évolution créatrice} p. 194 (-rattacher % PRINTED AS IS in the Danish
la vie soit à la conscience, soit à quelque chose qui y ressemble); p. 279 (la vie est d'ordre psychologique). --- Cf.\ on vitalism my Psykologi II, 3 and \textit{Mindre Arbejder.} First Series p. 40--50. --- Bergson has predecessors in French philosophy, especially \emph{Félix Ravaisson}. See the latter's \textit{La philosophie en France au 19\textsuperscript{e} siècle} p. 246. (On the earlier discussion of the question in French philosophy, see the same work p. 168--182).

\notelabel{37.} (p. 42.) On the romantic concept of evolution see \textit{Den nyere Filosofis Historie.}\textsuperscript{2} II p. 159--162. --- \textit{Mindre Arbejder.} Third Series. p. 205--207.

\notelabel{38.} (p. 44.) Cf., besides my \textit{Psykologi} (ch.\ IV and VII), my \textit{Mindre Arbejder.} Third Series. p. 28--53. --- Bergson's position on the problem of the will I have already discussed in \textit{Den menneskelige Tanke.} p. 291--295. --- Bergson's treatment of the problem is found in \textit{Les Données immédiates de la Conscience.} Chap.\ 3. De l'organisation des états de conscience: la liberté. The very title of this section is characteristic of his view.

\notelabel{39.} (p. 45.) What Bergson calls le moi fondamental or le moi d'en bas corresponds entirely to what I call ``the real self.'' Cf.\ \textit{Psykologi}\textsuperscript{6} p. 179--181; 354; 409--411.

\notelabel{40.} (p. 50.) Cf.\ my more detailed treatment in my \textit{Psykologi} VI E, 9.

\notelabel{41.} (p. 53.) \textit{L'intuition philosophique.} (Revue de Métaphysique et de Morale. 1911. p. 813).

\notelabel{42.} (p. 54.) \textit{Bulletin de la Société Française de Philosophie.} VIII p. 341. \textit{Introduction à la Métaphysique.} (Revue de Métaphysique et de Morale. 1903). p. 25 (quoted in ``Den menneskelige Tanke'' p. 318). In several places in \textit{Évolution créatrice} the inference by analogy is evident, see e.g.\ p. 190 f.; 250; 278--281; 387. --- In \textit{La Perception du Changement} p. 7--9 Bergson reproaches the philosophers with building their systems on a single part or side of existence, and he thinks he avoids this defect by immersing himself in our immediate apprehension itself, that from which all the conflicting concepts finally spring.

\notelabel{43.} (p. 54.) See more closely on analogies and their significance \textit{Den menneskelige Tanke.} p. 181--186; 310--315.

% --- p. 70 ---
\opage{70}\notelabel{44.} (p. 55) \textit{Évolution créatrice.} p. 270; 323. --- The definition of creation as growth is found on p. 262. % PRINTED AS IS in the Danish: "(p. 55)" without its stop

\notelabel{45.} (p. 57.) See on pragmatism \textit{Den menneskelige Tanke.} p. 270--272.

\notelabel{46.} (p. 57.) \textit{Évolution créatrice.} p. 290 f. --- \textit{Introduction à la Métaphysique.} p. 19 f.

\notelabel{47.} (p. 58.) \textit{La Perception du Changement} p. 36.

\notelabel{48.} (p. 59.) See more closely on this point my Helsingfors lectures \textit{Personlighetsprincipen i Filosofin.} (Stockholm 1911) p. 7--38, and the essay \textit{Aandelig Kultur} (Mindre Arbejder. Third Series) p. 4--6.

\notelabel{49.} (p. 59.) \textit{Évolution créatrice} p. 280.

\notelabel{50.} (p. 59.) \textit{Évolution créatrice} p. 192. --- \textit{La Perception du Changement.} p. 9--13; 36. --- \textit{L'Intuition Philosophique} (Revue de Métaphysique et de Morale 1912). p. 827.

\notelabel{51.} (p. 60.) More on this in my \textit{Religionsfilosofi} \S{} 44; 115--116; 131.

\notelabel{52.} (p. 61.) I know this discussion between Bergson and the Jesuit theologians only through \emph{Marcel Hébert's} essay: ``Henri Bergson et son affirmation de l'existence de Dieu.'' (Coenobium. Rivista internationale di liberi studi. 1912).

\end{small}


\end{document}
